\documentclass[11pt]{article}
\usepackage[T1]{fontenc}
\usepackage[utf8]{inputenc}
\usepackage{lmodern}
\usepackage{microtype}
\usepackage[margin=1in]{geometry}
\usepackage{amsmath,amssymb,amsthm,mathtools}
\usepackage{booktabs,longtable,array,tabularx}
\usepackage{enumitem}
\usepackage{xcolor}
\usepackage{listings}
\usepackage{xurl}
\usepackage[hidelinks]{hyperref}
\usepackage[nameinlink,noabbrev]{cleveref}
\usepackage{fancyhdr}
\usepackage{lastpage}

\hypersetup{
  pdftitle={A Unified Finitary-Relational Framework for Mathematics in iPRA},
  pdfauthor={OpenAI, prepared for Vladimir Reshetnikov},
  pdfsubject={Intuitionistic primitive recursive arithmetic, finite certificates, Lean formalization plan}
}

\pagestyle{fancy}
\fancyhf{}
\lhead{Finitary-relational mathematics in iPRA}
\rhead{15 November 2026}
\cfoot{\thepage\ of \pageref{LastPage}}
\setlength{\headheight}{14pt}
\setlength{\emergencystretch}{2em}

\newtheorem{theorem}{Theorem}[section]
\newtheorem{lemma}[theorem]{Lemma}
\newtheorem{corollary}[theorem]{Corollary}
\newtheorem{proposition}[theorem]{Proposition}
\newtheorem{definition}[theorem]{Definition}
\newtheorem{principle}[theorem]{Principle}
\theoremstyle{remark}
\newtheorem{remark}[theorem]{Remark}
\newtheorem{example}[theorem]{Example}
\newtheorem{status}[theorem]{Status}

\newcommand{\N}{\mathbb N}
\newcommand{\Z}{\mathbb Z}
\newcommand{\Q}{\mathbb Q}
\newcommand{\R}{\mathbb R}
\newcommand{\B}{\{0,1\}}
\newcommand{\ipra}{\mathsf{iPRA}}
\newcommand{\PRAAx}{\texttt{PRAAx}}
\newcommand{\Prf}{\texttt{Prf}}
\newcommand{\NPrf}{\texttt{NPrf}}
\newcommand{\PR}{\texttt{PR}}
\newcommand{\PTm}{\texttt{PTm}}
\newcommand{\Fm}{\texttt{Fm}}
\newcommand{\FRP}{\mathsf{FRP}}
\newcommand{\True}{\mathsf{true}}
\newcommand{\False}{\mathsf{false}}
\newcommand{\angles}[1]{\langle #1\rangle}
\newcommand{\code}[1]{\texttt{#1}}
\newcommand{\Close}{\mathsf{Close}}
\newcommand{\Good}{\mathsf{Good}}
\newcommand{\Compat}{\mathsf{Compat}}
\newcommand{\Viable}{\mathsf{Viable}}
\newcommand{\Obs}{\mathsf{Obs}}
\newcommand{\Dom}{\mathsf{Dom}}
\newcommand{\TraceOK}{\mathsf{TraceOK}}
\newcommand{\All}{\mathsf{All}}
\newcommand{\Any}{\mathsf{Any}}
\newcommand{\Search}{\mathsf{Search}}
\newcommand{\Run}{\mathsf{Run}}
\newcommand{\Inv}{\mathsf{Inv}}
\newcommand{\Accept}{\mathsf{Accept}}
\newcommand{\GoodTable}{\mathsf{GoodTable}}
\newcommand{\dist}{\mathsf{dist}}
\newcommand{\FullV}{\mathsf{FullV}}
\newcommand{\Reach}{\mathsf{Reach}}
\newcommand{\Tail}{\mathsf{Tail}}
\newcommand{\Next}{\mathsf{Next}}
\newcommand{\defeq}{\mathrel{:=}}
\newcommand{\dotminus}{\mathbin{\dot{-}}}

\lstdefinestyle{lean}{
  basicstyle=\ttfamily\small,
  columns=fullflexible,
  keepspaces=true,
  frame=single,
  breaklines=true,
  showstringspaces=false,
  keywordstyle=\bfseries,
  commentstyle=\itshape
}
\lstdefinestyle{python}{
  basicstyle=\ttfamily\small,
  columns=fullflexible,
  keepspaces=true,
  frame=single,
  breaklines=true,
  showstringspaces=false,
  language=Python
}

\title{\textbf{A Unified Finitary-Relational Framework\\for Embedding Mathematics into Intuitionistic Primitive Recursive Arithmetic}}
\author{Prepared for Vladimir Reshetnikov}
\date{15 November 2026\\Repository snapshot \code{b88e843c31b50dfe461bfbd0380166ca422cd3ab}}

\begin{document}
\maketitle

\begin{abstract}
The ProveIt repository contains a kernel-checked presentation of intuitionistic
primitive recursive arithmetic and a large collection of research packets that
translate classical mathematics into finite, quantitative forms.  The packets
repeatedly encounter the same boundaries: bounded induction is confused with
induction over completed infinite objects; ambient Lean theorems are confused
with object-language derivations; exact selectors are requested where only
multivalued approximate data are constructively justified; compactness is used
without a finite net or modulus; and deep finiteness theorems are used without
a height bound or certificate producer.

This report develops a single framework addressing these issues.  Its first
layer is a proof-producing finite-trace compiler: fixed primitive-recursive
algorithms, Boolean folds, bounded searches, and certificate verifiers are
translated into actual \(\Prf\,\PRAAx\) derivations by quantifier-free induction.
Its second layer is a finitary relational presentation of mathematical objects.
Instead of forcing a globally chosen representative, a theorem may return a
finite family of compatible approximate answers.  A backward viability dynamic
program amalgamates finite witnesses.  A new shrinking-viability theorem gives
a primitive-recursive coherent name when a quantitative determinacy modulus is
available; without such a modulus the correct output remains relational.

The report gives complete human proofs of the core compiler, Booleanization,
composition, witness-amalgamation, shrinking-viability, guarded extensionality,
finite-net transfer, and height-bound/list equivalence theorems.  It then shows
how these tools reorganize analysis, algebra, geometry, topology, number theory,
Diophantine geometry, and probability.  Exact selection principles known to
imply LLPO/LPO, universal primitive-recursive evaluation, and search cutoffs
from mere cardinality bounds are explicitly retained as barriers rather than
hidden assumptions.  The accompanying Lean files give dependency-ordered,
source-accurate blueprints against ProveIt's \code{ConciseIPRA} API; the generic
object-language proofs are not claimed kernel-checked in this archive.
\end{abstract}

\tableofcontents
\newpage

\section{Executive summary and exact scope}

\subsection{The problem exposed by the repository}

The repository separates two very different achievements.  The directory
\code{ConciseIPRA} contains a Lean formalization of ordinary first-order
intuitionistic primitive recursive arithmetic and its bi-interpretation with a
small native scalar calculus.  In particular, the repository provides the
object theory \(\PRAAx\), derivations \(\Prf\), quantifier-free induction
\code{PRA.qfInd}, and the transfer theorem \code{pra\_iff\_native}.  By contrast,
the theorem packets under \code{docs/proofs} are research reports: most contain
mathematical arguments, executable finite checks, and Lean plans, but none is a
completed derivation of its advertised theorem in \(\Prf\,\PRAAx\).

The recurring difficulty is therefore not merely ``how to code real numbers''
or ``how to avoid excluded middle.''  It is how to organize a proof so that
all of the following are simultaneously explicit:

\begin{enumerate}[label=(\arabic*)]
\item the finite information a construction reads;
\item the primitive-recursive algorithm or certificate checker operating on
      that information;
\item a quantifier-free invariant suitable for \(\ipra\)'s induction;
\item the modulus, separation margin, height bound, finite net, or covering
      data that makes an apparently infinitary step finite;
\item the exact object-language formula and derivation to be transported to
      the native calculus;
\item the boundary between a completed verifier and an unconstructed producer.
\end{enumerate}

The existing foundation packets already contain many of these ingredients.
The present report unifies them and adds two missing interfaces: a syntactic
proof-producing compiler and a relational refinement layer for non-unique
approximate witnesses.

\subsection{The proposed architecture}

The framework has five layers.

\paragraph{Layer 1: fixed numerical presentations.}
Every object family is represented by natural-number parameters and fixed
primitive-recursive evaluators.  A represented object is not an arbitrary
program supplied at run time.  If an external computation is not represented
by a fixed evaluator, its finite trace is supplied as a certificate and checked
by a fixed primitive-recursive verifier.  This distinction avoids the false
idea of a total universal primitive-recursive evaluator.

\paragraph{Layer 2: Boolean finite contracts.}
All bounded quantifiers, list checks, finite sums, maxima, residual tests, and
search predicates are compiled to primitive-recursive Boolean functions.
Global name hypotheses remain ordinary assumptions in the context; induction
uses only equations saying that a Boolean accumulator equals one.  Precision is
kept free during the finite proof and universally quantified only afterward.

\paragraph{Layer 3: proof-producing traces.}
A bounded computation is presented by fixed functions
\(\mathsf{init},\mathsf{step},\mathsf{bound},\mathsf{inv},\mathsf{accept}\).
Given object-language derivations of the base, local step, and final soundness
formulas, a generic theorem constructs an object-language derivation of the
whole run.  The construction uses \code{PRA.qfInd}, the defining equations of
primitive recursion, equality substitution, and ordinary natural deduction.
Thus the result is not merely an ambient theorem about a Lean function.

\paragraph{Layer 4: finitary relational presentations.}
A challenge may admit several correct approximate answers.  The framework
therefore permits a finite relation \(\Good(p,k,a)\) rather than insisting on a
single canonical \(a\).  A compatibility relation connects answers at adjacent
precision levels.  Bounded backward dynamic programming determines which
answers can be extended to a finite horizon and constructs a compatible finite
chain whenever one exists.  This is enough for all fixed-precision theorems and
avoids a hidden choice principle.

\paragraph{Layer 5: determinization only from quantitative data.}
When a downstream theorem genuinely needs one coherent name, it must consume a
quantitative determinacy hypothesis.  The shrinking-viability theorem assumes
that all candidates viable to a stated horizon lie in a small observational
ball.  The least viable candidate then defines a primitive-recursive name, and
the shrinking condition proves regularity.  Root separation, strong convexity,
isolating intervals, contraction estimates, or projective gaps are instances
of this hypothesis.  Qualitative uniqueness alone is not silently converted
into a modulus.

\subsection{Principal theorems}

The main general results proved in this report are:

\begin{enumerate}[label=\textbf{T\arabic*.},leftmargin=4em]
\item \textbf{Bounded Booleanization.} Every fixed bounded formula over
      primitive-recursive terms has a Boolean evaluator whose correctness is
      provable in \(\ipra\).
\item \textbf{Finite-trace compiler.} Embedded base, step, and final
      derivations yield an embedded theorem about the complete bounded run.
\item \textbf{Contract composition.} Proof-producing transducers compose,
      including adaptive finite query trees.
\item \textbf{Finite witness amalgamation.} A bounded backward viability
      computation constructs a compatible chain through finite candidate
      levels whenever such a chain exists.
\item \textbf{Shrinking viability.} Finite candidate relations plus a
      quantitative shrinking horizon define a coherent primitive-recursive
      approximate name.
\item \textbf{Guarded extensionality.} Operations with precision transformers
      descend to observational setoids, while partial operations are totalized
      and guarded by explicit domain or apartness certificates.
\item \textbf{Finite-net transfer.} A modulus and a finite net turn local
      rational checks into uniform conclusions on a represented compact set.
\item \textbf{Height bound/list equivalence.} For a decidable coded predicate,
      an explicit height bound and a complete finite list are constructively
      interconvertible; a mere cardinality bound is not enough.
\end{enumerate}

\subsection{What is and is not resolved}

The framework resolves several obstacles at the level at which they are
actually logical or architectural.  It gives a reusable path from finite
algorithms to \(\Prf\,\PRAAx\), a correct treatment of all-precision hypotheses,
a generic finite-search library, and a non-choice-based interface for
non-unique approximations.  It also gives standard positive replacements for
exact sign, root, axis, pivot, and branch selection.

It does not manufacture missing mathematics.  In particular, it does not
produce the unknown or ineffective height bounds needed for general Faltings
or Siegel statements, rational covering representatives for arbitrary
Mordell--Weil descent, or all constants in the strongest prime-distribution
claims.  Instead, those results are factored into a uniform finite verifier and
an explicit producer obligation.  This is progress: the verified downstream
arithmetic can be completed once a producer is supplied, and no search cutoff
is smuggled in through a qualitative theorem.

The included Lean files are detailed blueprints against the current repository
API.  They contain theorem statements, formula constructors, proof dependency
order, and exact points where \code{PRA.qfInd}, \code{Prf.subst},
\code{Prf.trans}, and \code{pra\_iff\_native} are used.  Lean was not available
in the artifact environment, so this archive does not claim that those files
have been elaborated or kernel-checked.

\section{Repository audit and diagnosis}

\subsection{Pinned source snapshot}

This report was prepared against the ProveIt main-branch snapshot
\[
\code{b88e843c31b50dfe461bfbd0380166ca422cd3ab}.
\]
At that snapshot, \code{Logic/PrimitiveRecursiveArithmetic/README.md}
distinguishes the checked \code{ConciseIPRA} library from 58 delivered theorem
packets under \code{docs/proofs}.  The repository's own index warns that the
packets were filed as delivered, were not rerun during intake, and do not
contain completed \(\Prf\,\PRAAx\) derivations of their advertised results.
That warning is foundational to the present audit: a polished PDF, a passing
Python test, a type-correct specification, or an ambient Lean theorem is not
itself an iPRA proof.

The collection spans algebra, real and complex analysis, analytic number
theory, combinatorics, Diophantine number theory, geometry, probability, and
topology.  Seven foundation packets already propose finite-contract,
finite-observation, certificate, modulus, and evidence-indexed approaches.
Their strongest common conclusions are sound and are retained here.  Their
main common gap is also retained: the last step from a generic semantic
argument to a constructed derivation in the actual \code{ConciseIPRA} object
language is mostly planned rather than implemented.

\subsection{Repeated proof-patterns}

Across the packets, successful arguments have a remarkably stable shape.

\begin{enumerate}
\item A classical object is replaced by a numerical presentation: a finite
      coefficient list, a fixed Cauchy evaluator, a rational law, a finite
      grid, a coded matrix, a bounded trace, or a supplied certificate.
\item The desired statement is rewritten as a family of finite inequalities,
      coefficient identities, residual bounds, parity checks, or finite
      search assertions, indexed by a requested precision or bound.
\item Every non-total operation is guarded by a denominator, derivative,
      separation, pivot, or compactness margin.
\item A primitive-recursive calculation is justified by a local invariant.
\item The final mathematical statement is recovered by universal introduction
      over the previously free precision parameter or by an extensionality
      theorem for the chosen presentation.
\end{enumerate}

The same pattern appears in the constructive prime number theorem, finite
Euler products, polynomial factorization, exact coordinate geometry,
finite-grid topology, rational probability laws, and bounded Diophantine
enumeration.  This is evidence that a common framework is appropriate.

\subsection{Repeated failure-modes}

The difficult packets also repeat a small number of failure-modes.

\paragraph{Ambient proof versus embedded derivation.}
Many Lean files prove correct facts about \code{Nat}, \code{Rat}, arrays, or
finite lists.  These are useful semantic prototypes, but they do not
construct a term of type \(\Gamma\vdash[\PRAAx]\varphi\).  The bridge theorem
\code{pra\_iff\_native} transports an already constructed iPRA derivation; it
does not synthesize that derivation from an ambient theorem.

\paragraph{Induction over an infinitary interpretation.}
A represented-real contract may be written as an assertion quantified over
all precisions.  Placing that assertion inside the induction predicate makes
the predicate non-quantifier-free and creates the impression that iPRA is too
weak.  The correct proof fixes the requested precision, instantiates only the
finitely many name facts read by the computation, and inducts on a Boolean
accumulator.

\paragraph{Exact choice at unstable points.}
Root ordering at collisions, a unit rotation axis at the identity, exact sign
of an arbitrary real, a unique minimizer without a modulus, or a Puiseux
branch from arbitrary Cauchy coefficients can be discontinuous.  Several
packets correctly identify LLPO/LPO or computability obstructions.  The right
repair is a relational or approximate output, or an exact output under an
explicit positive gap.

\paragraph{Compactness without data.}
Statements such as ``take a finite subcover,'' ``choose a convergent
subsequence,'' or ``minimize on a compact set'' have no direct computational
content.  In the successful packets they are replaced by a finite net, a
modulus of uniform continuity, a finite cover certificate, or a nested family
of isolating regions.

\paragraph{Finiteness without a cutoff.}
A proof that a set has at most \(M\) elements does not identify where those
elements occur.  General Faltings and Siegel packets isolate precisely this
gap: bounded enumeration is straightforward once a height bound is supplied,
but cardinality alone does not give a stopping rule.  The same issue appears
in rational covering representatives for Mordell--Weil descent.

\paragraph{Arbitrary program codes.}
A theorem schema may be uniform in the numerical parameters of one fixed
primitive-recursive evaluator.  It cannot quantify over arbitrary program
codes and then invoke a total primitive-recursive universal evaluator.  A
fixed evaluator or a finite trace verifier is required.

\subsection{What the earlier foundation packets contribute}

The existing foundation work can be organized as follows.

\begin{longtable}{@{}p{0.24\textwidth}p{0.68\textwidth}@{}}
\toprule
Packet idea & Contribution retained here \\
\midrule
\endhead
Certificate-and-modulus framework & Bounded search, certificate composition,
regular-name arithmetic, apartness-conditioned inverse, and a clear
producer/verifier split.\\
Finite contracts & Booleanization, guarded induction, fixed evaluators,
finite input tables, stable windows, local analytic branches, and explicit
power-gap certificates.\\
Finite observations & Separation of evaluation and verification demands,
backward precision schedules, finite-net infima, and examples of checked
finite-table Lean proofs.\\
Certified Finitary Presentations & A clean trace-system master theorem,
finite query tables, fixed-precision discipline, domain-specific resolution
matrix, and explicit computability/constructive barriers.\\
Evidence-indexed finite mathematics & Relational outputs and backward finite
viability, avoiding an unjustified globally chosen approximate witness.\\
Uniform certificates & Finite-net lifting, sign-bracket root names, and
precision-contract composition.\\
Targeted embedding framework & Simple analytic branches, simultaneous finite
window stability, finite algebra certificates, and a hyperelliptic fiber
count conditional on a missing covering theorem.\\
\bottomrule
\end{longtable}

The present report does not claim priority over those ideas.  Its contribution
is to combine them into one proof architecture, strengthen the relational
component into a general shrinking-viability theorem, and state the
object-language compiler in a form whose inputs and output are actual
\(\Prf\,\PRAAx\) derivations.

\subsection{A three-way status discipline}

Every claim in a future packet should be classified as one of the following.

\begin{description}[leftmargin=3.4cm,style=nextline]
\item[Embedded theorem.] A checked object-language derivation, optionally exported to the native
      calculus.
\item[Semantic theorem.] A complete human proof or ambient Lean proof about
      the coded structures, accompanied by an exact plan for the embedded
      proof, but not yet a \(\Prf\) term.
\item[Target or conditional theorem.] A well-formed proposition whose proof
      depends on a named producer, modulus, height bound, covering list, or
      analytic certificate not supplied by the packet.
\end{description}

This report's core mathematical arguments are semantic theorems with detailed
embedded proof plans.  The accompanying Python checks are implementation
examples.  No line of the report treats those checks as a replacement for a
universal proof.

\section{Logical base and the meaning of an embedded proof}

\subsection{The repository's ordinary presentation of iPRA}

The checked backend uses the following objects.

\begin{itemize}
\item \(\PR\,k\) is the inductively generated type of primitive-recursive
      function symbols of arity \(k\), with zero, successor, projections,
      composition, and primitive recursion.
\item \(\PTm\,n\) is the first-order term language with \(n\) free variables,
      formed by variables and applications of \(\PR\) symbols.
\item \(\PRAAx\) contains the defining equations for every primitive-recursive
      symbol, successor separation and injectivity, and induction for
      quantifier-free formulas.
\item \(\Prf\,\PRAAx\,\Gamma\,\varphi\), written
      \(\Gamma\vdash[\PRAAx]\varphi\), is intuitionistic natural deduction
      with equality, substitution, and open axioms.
\end{itemize}

The crucial repository theorem is
\[
(\Gamma\vdash[\PRAAx]\varphi)
\quad\Longleftrightarrow\quad
\NPrf(\Gamma^{\mathrm{nat}},\varphi^{\mathrm{nat}}),
\]
implemented as \code{pra\_iff\_native}.  It is a faithful transport theorem,
not a reflection principle from arbitrary Lean propositions.

\subsection{Schemas versus single first-order theorems}

A fixed primitive-recursive evaluator \(a:\N^{r+1}\to\N\) can be represented
by one fixed \(\PR\) symbol.  A theorem may then be uniform in its ordinary
numerical parameter \(p\) and precision \(k\):
\[
  \forall p\,\forall k\;\Phi_a(p,k).
\]
This is a single first-order theorem for that fixed symbol.  A Lean theorem
polymorphic in a meta-level value \code{a : PR (r+1)} is instead a schema
generator: for each fixed symbol it produces a first-order formula and proof.
The distinction is benign and should be explicit.

What is not legitimate is to quantify over arbitrary program codes \(e\) and
pretend that there is a total primitive-recursive function
\(U(e,x)\) evaluating every primitive-recursive program.  The diagonal
argument in \cref{sec:barriers} rules this out.  A theorem about an external
algorithm must therefore use one of two interfaces:

\begin{enumerate}
\item its evaluator is fixed in the theorem's language; or
\item a finite computation trace is supplied and checked by a fixed verifier.
\end{enumerate}

\subsection{Definitional extensions are harmless when fixed}

The usual mathematical presentation of iPRA has a function symbol for every
primitive-recursive function.  The repository instead represents these
symbols by syntax trees in \(\PR\), with defining equations in \(\PRAAx\).
This means that introducing a named helper such as pairing, list lookup,
integer addition, rational comparison, bounded maximum, or a particular trace
runner is not an additional axiom.  One defines a fixed \(\PR\) term and uses
its generated defining equations.

For proof engineering it is useful to distinguish three levels.

\begin{description}[leftmargin=3.2cm,style=nextline]
\item[Meta-level definition.] A Lean function constructs a \(\PR\) syntax
      tree.  This is not yet a theorem.
\item[Object term.] Applying the resulting symbol gives a \(\PTm\) term.
\item[Object proof.] Theorems such as \code{PRA.comp}, \code{PRA.precZero},
      and \code{PRA.precSucc} expose the defining equations as
      \(\Prf\,\PRAAx\) derivations.
\end{description}

A formalization should avoid duplicating encodings.  The same Boolean, pair,
list, signed-integer, and positive-denominator rational codes should be reused
across all packets, because every additional representation creates a new
extensionality and conversion burden.

\subsection{Quantifier-free induction with quantified context}

A frequent misunderstanding is that quantifier-free induction forbids the use
of represented-object hypotheses containing quantifiers.  It does not.  The
restriction applies to the induction formula, not to the surrounding context.
For example, suppose the context contains a regularity hypothesis
\[
  H(p):\equiv \forall k\,\forall \ell\;
    \dist(a(p,k),a(p,\ell))\le 2^{-k}+2^{-\ell}.
\]
Inside the successor step one may instantiate \(H(p)\) at primitive-recursive
terms computed from the current finite state.  The induction predicate itself
can remain the equation
\[
  \mathsf{accumulator}(i,p)=1.
\]
Thus a global quantified contract may be consumed stage by stage while the
formula on which induction acts is quantifier free.

This observation is the logical basis of the finite query-table discipline.
A computation declares or records every sample it reads.  The global contract
is instantiated only at those samples, and the resulting finite facts are
folded into a Boolean table.  The induction invariant says that the prefix of
the table already constructed is valid.

\subsection{The proof-producing interface}

The correct reusable abstraction is not an ambient proposition asserting that
an algorithm is sound.  It is a structure whose soundness field is already an
object-language derivation.

At a schematic level, a proof-producing transducer from a contract \(A\) to a
contract \(B\) consists of:
\[
  \text{an output term }F(\bar x),
  \qquad
  A(\bar x)\ \vdash[\PRAAx]\ B(\bar x,F(\bar x)).
\]
Composition is object-language cut.  If
\[
 A\vdash B(F) \quad\text{and}\quad B(y)\vdash C(G(y)),
\]
then substitution followed by \code{Prf.trans} gives
\[
 A\vdash C(G(F)).
\]
This simple formulation prevents a whole class of category mistakes.  The
semantic development may be large, but its reusable interface ends in a term
of the repository's derivation type.

\subsection{What the generic framework may assume}

The core theorems below assume a small arithmetic library whose existence is
standard and whose completion is a finite formalization task:

\begin{enumerate}
\item primitive-recursive pairing and finite-list coding, with projection and
      length equations;
\item Boolean values \(0,1\), Boolean conjunction/disjunction/negation, and
      normalization lemmas;
\item zero test, equality test, order, truncated subtraction, and bounded
      search;
\item signed integers and rationals represented by normalized or guarded
      codes; and
\item dyadic radii and elementary error-budget arithmetic.
\end{enumerate}

Nothing in the master trace theorem depends on the details of these codes.
For a concrete Lean implementation, however, selecting one canonical library
before building domain-specific modules is the first milestone.

\section{Bounded Booleanization and finite search}
\label{sec:booleanization}

\subsection{Boolean predicates}

A numerical predicate \(P(\bar x)\) is represented by a primitive-recursive
function \(b_P(\bar x)\in\B\).  We write
\[
  [P](\bar x) \quad\text{for the equation}\quad b_P(\bar x)=1.
\]
The Boolean operations are fixed primitive-recursive functions satisfying the
truth tables for \(0\) and \(1\).  Their correctness is proved by the two
numerical cases.  Equality and order are also represented by Boolean functions;
for example, \(x\le y\) can be tested by whether truncated subtraction
\(x\dotminus y\) is zero.

A \emph{bounded formula} is built from atomic equations between
primitive-recursive terms, propositional connectives, and quantifiers
\(\forall i<B(\bar x)\) and \(\exists i<B(\bar x)\).  The bound and all terms
are fixed primitive-recursive functions.

\subsection{Finite folds}

For a Boolean predicate \(b(i,\bar x)\), define
\begin{align*}
 \All_b(0,\bar x)&=1,\\
 \All_b(n+1,\bar x)&=\All_b(n,\bar x)\wedge b(n,\bar x),\\[3pt]
 \Any_b(0,\bar x)&=0,\\
 \Any_b(n+1,\bar x)&=\Any_b(n,\bar x)\vee b(n,\bar x).
\end{align*}
Both are primitive recursive.  Define a first-witness search
\(\Search_b(n,\bar x)\) that returns the least \(i<n\) with
\(b(i,\bar x)=1\), and returns \(n\) when no such witness exists.  It can be
computed by a primitive recursion carrying the current candidate and a found
bit.

\begin{theorem}[Correctness of Boolean folds]
\label{thm:boolean-fold}
For every fixed Boolean predicate \(b\), iPRA proves:
\begin{align*}
 \All_b(n,\bar x)=1
  &\leftrightarrow \forall i\,(i<n\to b(i,\bar x)=1),\\
 \Any_b(n,\bar x)=1
  &\leftrightarrow \exists i\,(i<n\wedge b(i,\bar x)=1),
\end{align*}
and, under the existential premise,
\[
 \Search_b(n,\bar x)<n,
 \qquad b(\Search_b(n,\bar x),\bar x)=1,
\]
with no smaller witness.
\end{theorem}

\begin{proof}
We spell out the proof because it illustrates the permitted use of
quantifier-free induction.

For the forward direction of the universal fold, keep \(i\) free and induct on
\(n\) with the quantifier-free formula
\[
  i<n\wedge \All_b(n,\bar x)=1\;\to\;b(i,\bar x)=1.
\]
The relation \(i<n\) and the Boolean equation are atomic equations between
primitive-recursive terms.  At zero the antecedent is false.  At a successor,
if \(i<n+1\), numerical case analysis gives either \(i<n\), where the induction
hypothesis applies to the prefix fold, or \(i=n\), where the last Boolean
factor gives \(b(n,\bar x)=1\).  Universal introduction over \(i\) is performed
after the induction.

For the reverse direction, place
\(H:=\forall i(i<n\to b(i,\bar x)=1)\) in the context and induct on the fold
length \(m\le n\) with the quantifier-free formula
\(\All_b(m,\bar x)=1\).  The step specializes \(H\) to \(m\), uses the
recursion equation of \(\All_b\), and the Boolean truth table for conjunction.
Equivalently, one may prove the open theorem with \(n\) as induction variable,
keeping \(\forall i(i<n\to\cdots)\) outside the induction predicate and using
its successor decomposition in the context.

The existential fold is dual.  For the forward direction, induction on \(n\)
examines whether the last disjunct is true; the result is either the last index
or a witness obtained from the prefix induction hypothesis.  For the reverse
direction, a supplied witness \(i<n\) is propagated through the fold by
quantifier-free induction on the difference between \(i+1\) and \(n\).

For first-witness search, the invariant states that either the found bit is
zero and every earlier predicate value is zero, or the found bit is one, the
stored index is below the current length, its predicate value is one, and all
smaller indices fail.  Every clause is a Boolean equation or bounded
comparison, hence quantifier free.  At termination, an existential premise
excludes the no-witness branch.  The minimality clause is read directly from
the invariant.
\end{proof}

\subsection{Booleanization theorem}

\begin{theorem}[Bounded Booleanization]
\label{thm:booleanization}
For every fixed bounded formula \(\theta(\bar x)\), one can construct a
primitive-recursive Boolean function \(b_\theta\) such that iPRA proves
\[
  b_\theta(\bar x)=1\quad\leftrightarrow\quad\theta(\bar x).
\]
The construction is effective in the syntax of \(\theta\).
\end{theorem}

\begin{proof}
Proceed by structural recursion on \(\theta\).  Atomic equations use the
Boolean equality test.  Conjunction, disjunction, and implication use the
corresponding Boolean truth functions; intuitionistic logic proves the finite
truth-table equivalences because each input is known to equal zero or one.
Negation is implication to the designated false equation.

For \(\forall i<B(\bar x)\,\psi(i,\bar x)\), apply the induction hypothesis to
obtain \(b_\psi\), and define
\[
 b_\theta(\bar x)=\All_{b_\psi}(B(\bar x),\bar x).
\]
The result follows from \cref{thm:boolean-fold}.  The existential case uses
\(\Any\).  All constructions use composition and primitive recursion, so the
resulting function is represented by a fixed \(\PR\) symbol.  The proof is
uniformly assembled from the component equivalences by substitution and
natural-deduction rules.
\end{proof}

\subsection{Bounded minimization}

Let \(c(i,\bar x)\) be an integer or rational cost with a decidable comparison,
and let \(C(\bar x)>0\) be a finite candidate bound.  Define
\(\operatorname{argmin}_{i<C}c(i,\bar x)\) by primitive recursion, retaining
the least index among equal minima.

\begin{lemma}[Finite minimization]
\label{lem:finite-min}
I P R A proves that the returned index is below \(C\) and has cost no larger
than every candidate below \(C\).  If a Boolean admissibility predicate is
included and at least one candidate is admissible, the same holds among
admissible candidates.
\end{lemma}

\begin{proof}
The accumulator stores the current best index, a found bit, and optionally the
current best cost.  The quantifier-free invariant states that the best index is
admissible when the found bit is one and that every previously inspected
admissible candidate has cost at least the stored cost; tie-breaking adds that
no smaller equal-cost candidate was skipped.  The successor step is one
comparison and a Boolean case split.  The existential admissibility premise
eliminates the no-candidate final state.
\end{proof}

This lemma is the computational core of finite-net residual minimization,
finite maximum and minimum, canonical certificate selection, and the backward
viability algorithms below.

\subsection{Bounded descent}

A common proof has a measure \(m(s)\in\N\) that strictly decreases whenever a
state is not terminal.  If the initial measure is bounded by \(M(u)\), then the
state must terminate within \(M(u)\) steps.

\begin{lemma}[Primitive-recursive descent]
\label{lem:descent}
Suppose fixed primitive-recursive functions \(\mathsf{step}\),
\(\mathsf{terminal}\), and \(m\) satisfy, in iPRA,
\[
 \mathsf{terminal}(s)=0\to m(\mathsf{step}(s))<m(s).
\]
Then the bounded run of length \(m(s)+1\) reaches a terminal state.
\end{lemma}

\begin{proof}
Induct on the step index with the quantifier-free invariant
\[
 \mathsf{terminal}(s_i)=1
 \quad\vee\quad m(s_i)+i\le m(s_0).
\]
If the state is already terminal, keep it fixed.  Otherwise the strict measure
decrease proves the second disjunct at the next stage.  At stage
\(m(s_0)+1\), the nonterminal disjunct would imply a negative natural measure,
so only terminality remains.
\end{proof}

The point is not that all mathematical descent arguments are easy, but that
once a numerical measure and local decrease theorem are supplied, termination
is a routine quantifier-free induction.

\section{The proof-producing finite-trace compiler}
\label{sec:trace-compiler}

\subsection{Trace systems}

All run-time data are paired into one natural number \(u\): original inputs,
requested precision, numerical bounds, and externally supplied certificate
codes.  A \emph{trace system} consists of fixed primitive-recursive functions
\[
\begin{array}{lll}
 \mathsf{good}(u), & \mathsf{init}(u), & \mathsf{bound}(u),\\
 \mathsf{step}(i,s,u), & \mathsf{inv}(i,s,u), &
 \mathsf{accept}(s,u),
\end{array}
\]
where the Boolean-valued functions are normalized to \(0\) or \(1\).  Its run
is the primitive recursion
\begin{align*}
 \Run(0,u)&=\mathsf{init}(u),\\
 \Run(i+1,u)&=\mathsf{step}(i,\Run(i,u),u).
\end{align*}
In the repository's \(\PR\) syntax this is exactly
\(\code{PR.prec init step}\), with one parameter after the recursion argument.

The local proof obligations are open object-language formulas:
\begin{align}
 \mathsf{good}(u)=1 &\to
   \mathsf{inv}(0,\mathsf{init}(u),u)=1,
   \tag{B}\label{eq:trace-base}\\
 \mathsf{good}(u)=1\wedge\mathsf{inv}(i,s,u)=1 &\to
   \mathsf{inv}(i+1,\mathsf{step}(i,s,u),u)=1,
   \tag{S}\label{eq:trace-step}\\
 \mathsf{good}(u)=1\wedge
   \mathsf{inv}(\mathsf{bound}(u),s,u)=1 &\to
   \mathsf{accept}(s,u)=1.
   \tag{F}\label{eq:trace-final}
\end{align}
The essential requirement is not merely that these implications be true in
\(\N\).  A proof-producing trace package contains actual derivations
\[
 []\vdash[\PRAAx]\eqref{eq:trace-base},\qquad
 []\vdash[\PRAAx]\eqref{eq:trace-step},\qquad
 []\vdash[\PRAAx]\eqref{eq:trace-final},
\]
with the displayed free variables.  They may themselves be generated from
smaller proof-producing components.

\subsection{The compiler theorem}

\begin{theorem}[Embedded finite-trace reflection]
\label{thm:embedded-trace}
From object-language derivations of \eqref{eq:trace-base}--\eqref{eq:trace-final}
one can construct an object-language derivation
\[
 []\vdash[\PRAAx]
 \forall u\;\bigl(
  \mathsf{good}(u)=1\to
  \mathsf{accept}(\Run(\mathsf{bound}(u),u),u)=1
 \bigr).
\]
The construction is primitive recursive on the supplied proof terms.
\end{theorem}

\begin{proof}
We give the proof at the level of the repository's derivation constructors.
Keep \(u\) free and place \(G(u):\mathsf{good}(u)=1\) in the context.  Define
the induction formula with free variables \(i,u\):
\[
 A(i,u):\equiv
 \mathsf{inv}(i,\Run(i,u),u)=1.
\]
It is quantifier free because it is one equation between terms formed from
fixed primitive-recursive symbols.

For the base derivation, instantiate \eqref{eq:trace-base} at the current
\(u\), apply it to the context hypothesis \(G(u)\), and use the defining
zero equation
\[
 \Run(0,u)=\mathsf{init}(u)
\]
provided by \code{PRA.precZero}.  Equality substitution changes the resulting
invariant at \(\mathsf{init}(u)\) into \(A(0,u)\).

For the successor derivation, enter the induction context containing
\(A(i,u)\) and the shifted copy of \(G(u)\).  Instantiate
\eqref{eq:trace-step} with
\(s:=\Run(i,u)\).  Modus ponens gives
\[
 \mathsf{inv}(i+1,
   \mathsf{step}(i,\Run(i,u),u),u)=1.
\]
The successor defining equation supplied by \code{PRA.precSucc} is
\[
 \Run(i+1,u)=\mathsf{step}(i,\Run(i,u),u).
\]
Use \code{Prf.eqSubst1} (or simultaneous \code{Prf.eqSubst}) to rewrite the
state argument.  This produces \(A(i+1,u)\).

Apply \code{PRA.qfInd} to obtain, in context \(G(u)\),
\(\forall i\,A(i,u)\).  Eliminate the universal quantifier at the term
\(\mathsf{bound}(u)\).  Instantiate \eqref{eq:trace-final} with
\(s:=\Run(\mathsf{bound}(u),u)\) and apply it to \(G(u)\) and the specialized
invariant.  Discharge \(G(u)\) by implication introduction and close \(u\) by
universal introduction.

Every operation in this construction---substitution of terms into a supplied
derivation, weakening, implication elimination, equality substitution,
quantifier-free induction, and universal introduction---is a constructor or a
proved transformation in \code{ConciseIPRA.Deduction} and
\code{ConciseIPRA.PR}.  Consequently the construction can be implemented as a
Lean function producing a \(\Prf\) term, not merely as a metatheoretic
soundness argument.
\end{proof}

\begin{remark}[Why the theorem is nontrivial but modest]
The theorem does not assert that every true primitive-recursive statement is
provable in iPRA.  Its inputs already include proofs of the local arithmetic
obligations.  It states that bounded execution and invariant propagation add
no further logical strength: they are compiled by quantifier-free induction.
The mathematical burden is to design an invariant for which the local step is
actually derivable.
\end{remark}

\subsection{Finite query tables}

Suppose an input object is supplied through a fixed evaluator \(a(p,k)\) and a
global contract such as
\[
 H_a(p):\equiv\forall k\,\forall\ell\;
 R(a(p,k),a(p,\ell),k,\ell).
\]
A trace at requested precision \(n\) reads only finitely many samples.  Queries
may be adaptive: the next precision can depend on previous rational
comparisons.  Include in the trace state a query log and a Boolean accumulator
recording all instantiated contract facts.

Define \(\GoodTable(p,n,c)=1\) to be the finite conjunction of exactly the
regularity, domain, magnitude, and separation facts needed along the run coded
by \(c\).  If the algorithm computes its own query log, \(c\) may simply be the
run state.

\begin{lemma}[Finite instantiation]
\label{lem:finite-query}
For every fixed query generator, iPRA proves
\[
 H_a(p)\to \GoodTable(p,n,c)=1
\]
for the actual finite query table, provided each table field is an instance or
a finite algebraic consequence of \(H_a(p)\) and any other global hypotheses.
\end{lemma}

\begin{proof}
Induct over the finite query trace with a Boolean equation saying that all
entries recorded so far are valid.  At a successor stage, the next query
indices are primitive-recursive terms in the previous state.  Universal
elimination instantiates \(H_a(p)\) at those terms.  The resulting finite
arithmetic fact is inserted into the accumulator using a proved Boolean
conjunction lemma.  The induction formula contains only the accumulator
equation; the quantified hypothesis remains in the context.
\end{proof}

This is the exact repair for ``induction over all precisions.''  The requested
precision \(n\) is free, the run consumes a finite table, and only after the
fixed-\(n\) proof is complete is \(n\) universally introduced.

\begin{corollary}[Fix precision, then quantify]
\label{cor:fix-precision}
Suppose global hypotheses imply the finite input table for every free precision
\(n\), the trace laws are embedded derivations, and acceptance implies a finite
output contract \(\mathsf{Out}(p,n,A(p,n))\).  Then iPRA proves
\[
 H(p)\to\forall n\;\mathsf{Out}(p,n,A(p,n)).
\]
No induction formula contains \(\forall n\).
\end{corollary}

\subsection{Certificate mode}

A large computation need not be rerun by primitive recursion inside the final
theorem.  A certificate may contain the state sequence
\(s_0,\ldots,s_B\), local arithmetic witnesses, and the output.  A fixed
Boolean verifier checks
\begin{align*}
 s_0&=\mathsf{init}(u),\\
 s_{i+1}&=\mathsf{step}(i,s_i,u)\quad(i<B),\\
 \mathsf{inv}(i,s_i,u)&=1\quad(i\le B),\\
 \mathsf{accept}(s_B,u)&=1.
\end{align*}
Every clause is a bounded fold over the coded list.

\begin{theorem}[Certificate reflection]
\label{thm:certificate-reflection}
For a fixed certificate language and checker, one can construct an embedded
derivation
\[
 []\vdash[\PRAAx]\forall u\,\forall c\;
 \bigl(\TraceOK(u,c)=1\to
       \mathsf{Conclusion}(u,\mathsf{output}(c))\bigr),
\]
provided the local checker clauses have embedded soundness derivations.
\end{theorem}

\begin{proof}
Boolean-fold correctness exposes each checked list clause.  There are two
proof routes.  The direct route reads the final acceptance clause and applies
its embedded soundness theorem.  The reconstructive route proves by
quantifier-free induction on the list index that the certified state equals the
canonical run, then invokes \cref{thm:embedded-trace}.  The direct route is
usually smaller; the reconstructive route is useful when later theorems need
the canonical run equality.
\end{proof}

Certificate mode is appropriate for polynomial identities, isolating disks,
factorization traces, finite covers, high-precision analytic constants, and
large numerical inequalities.  A separate producer theorem may later show
that a certificate exists for a given input family.

\subsection{Sequential and network composition}

A proof-producing transducer \(F:X\to Y\) consists of a primitive-recursive
output algorithm, a finite input-demand function, and an embedded soundness
derivation.  If \(G:Y\to Z\) is another transducer, their composite first runs
\(F\), extracts the finite metadata needed to determine \(G\)'s demand, and
then runs \(G\).

\begin{theorem}[Composition]
\label{thm:trace-composition}
Proof-producing finite transducers are closed under sequential composition,
finite products, finite case splits with overlapping guards, and finite
acyclic networks.  The composite has primitive-recursive demand and bound
functions and an embedded soundness derivation.
\end{theorem}

\begin{proof}
For sequential composition, encode the state as a phase bit, a component state,
and frozen metadata.  Phase zero executes \(F\).  On its final step, the state
stores the certified output and all magnitude/domain data needed to initialize
\(G\).  Phase one executes \(G\).  The total bound is the sum of the two bounds
plus a switch step.

The invariant is a Boolean case split on the phase.  In phase zero it is
\(F\)'s invariant.  At the switch it includes \(F\)'s acceptance and \(G\)'s
input contract.  In phase one it is \(G\)'s invariant together with the frozen
metadata.  The local step proof is assembled from the component step proofs,
the switch soundness proof, and Boolean case analysis.  Apply
\cref{thm:embedded-trace}.  The final implication is \(G\)'s soundness after
substituting \(F\)'s certified output.

Products pair states and conjoin invariants.  Overlapping cases use a total
selector whose branch-correctness proof is supplied in each overlap.  A finite
acyclic network is topologically ordered and flattened into phases.  All proof
composition is substitution plus \code{Prf.trans} or ordinary implication
elimination.
\end{proof}

\subsection{The exact compiler boundary}

The trace compiler is a disciplined boundary between semantic mathematics and
formal arithmetic.

\begin{itemize}
\item Outside the boundary one proves local rational inequalities,
      coefficient identities, error budgets, and invariant preservation.
\item At the boundary each local theorem is expressed as a formula of
      \(\Fm\,\PTm\) and supplied as a \(\Prf\,\PRAAx\) derivation.
\item Inside the boundary generic code combines those derivations into the
      complete bounded-run theorem.
\item Only after that theorem exists is \code{pra\_iff\_native} applied.
\end{itemize}

This organization makes missing work visible.  A semantic Lean theorem can be
used as a guide for constructing a local derivation, but it cannot be inserted
as an unexamined premise.  Conversely, once the local obligations are embedded,
the generic control-flow proof should not be repeated theorem by theorem.

\section{Finitary relational presentations}
\label{sec:relational}

\subsection{Why single-valued names are sometimes the wrong interface}

Many constructive failures arise because a theorem is asked to return one
exactly selected representative when the mathematics only determines a finite
family of approximations.

\begin{itemize}
\item A polynomial with colliding roots has an unambiguous root multiset but no
      stable ordering of individual roots.
\item A real rotation near the identity has stable quaternion or approximate
      axis data, but no continuous choice of a unit axis at the identity.
\item A minimization problem may have several nearly minimizing grid points,
      even though the minimum value is stable.
\item Local analytic charts can overlap, while an exact chart selector at a
      boundary is discontinuous.
\item A finite proof may establish a compatible witness at every requested
      depth without supplying one globally chosen infinite branch.
\end{itemize}

Forcing a function in these situations either adds a separation hypothesis,
imports a choice principle, or makes the theorem false as a uniform
computational statement.  The framework therefore treats relations as a
first-class output type.

\subsection{Definition of a finitary relational presentation}

\begin{definition}[Finitary relational presentation]
A finitary relational presentation \(X\) consists of fixed
primitive-recursive data:
\begin{enumerate}
\item a parameter-domain Boolean \(\Dom_X(p)\);
\item a challenge-domain Boolean \(\mathsf{Chal}_X(p,q)\);
\item a finite answer bound \(B_X(p,q)\);
\item an answer verifier \(\Good_X(p,q,a)\in\B\), with the convention that
      valid answers satisfy \(a<B_X(p,q)\);
\item a refinement or compatibility verifier
      \(\Compat_X(p,q,a,q',a')\in\B\); and
\item optional observation and error decoders used to interpret an answer.
\end{enumerate}
All data are fixed symbols, not arbitrary run-time programs.
\end{definition}

The challenge \(q\) need not be a scalar precision.  It may encode a region,
derivative order, finite test function, matrix entry, truncation length, or a
bundle of such requests.  This makes the interface broad enough for analysis,
probability, topology, and algebra without pretending that every object is
merely a Cauchy sequence.

A parameter \(p\) is \emph{finitely inhabited to depth \(h\)} when every
challenge in the finite schedule up to \(h\) has an acceptable answer and the
answers can be connected by the required compatibility edges.  The exact
schedule is theorem-specific and primitive recursive.

\subsection{Functional presentations as a special case}

A usual fixed evaluator \(a(p,q)\) gives a relational presentation by setting
\[
 \Good_X(p,q,z)=1 \quad\Longleftrightarrow\quad z=a(p,q),
\]
and taking compatibility to be the desired regularity relation.  Thus no
existing functional development is lost.  The relational form is strictly
more flexible: it can represent finite root clusters, all valid pivots, all
near-minimizers, all cells intersecting a curve, or all charts whose guards
hold.

\subsection{Relational transducers}

\begin{definition}[Relational transducer]
A finitary relational transducer \(F:X\rightsquigarrow Y\) consists of:
\begin{enumerate}
\item a primitive-recursive finite query strategy.  Given an input parameter
      \(p\), output challenge \(r\), and current finite state, it either asks
      for an \(X\)-answer to a challenge \(q\) or halts;
\item a primitive-recursive output-answer bound and candidate generator;
\item a Boolean trace invariant and acceptance predicate; and
\item an embedded soundness derivation saying that any finite family of valid,
      mutually compatible input answers accepted by the trace yields a valid
      \(Y\)-answer.
\end{enumerate}
\end{definition}

The transducer need not choose among all valid input answers.  Its theorem is
relational: every accepted compatible input package leads to an acceptable
output package, or it produces a finite set of output candidates with a
verified nonemptiness and compatibility relation.

\subsection{Composition without hidden choice}

Suppose \(F:X\rightsquigarrow Y\) and \(G:Y\rightsquigarrow Z\).  A naive
composition would select one \(Y\)-answer for each query made by \(G\).  That
selection can destroy compatibility.  The relational composition instead
records a finite family of \(Y\)-candidates produced by \(F\), together with
their compatibility graph, and asks whether a compatible subfamily satisfies
\(G\)'s finite demand.  This is a finite constraint problem and is solved by
the viability dynamic program of \cref{sec:amalgamation}.

\begin{theorem}[Relational composition]
\label{thm:rel-composition}
The composition of two finitary relational transducers is a finitary
relational transducer.  Its verifier and finite candidate bounds are primitive
recursive.  Its soundness derivation is obtained by the embedded trace
compiler and the finite witness-amalgamation theorem.
\end{theorem}

\begin{proof}
Flatten the two traces into one phase-indexed trace as in
\cref{thm:trace-composition}.  Whenever \(G\) requests a \(Y\)-answer, invoke
\(F\)'s finite candidate generator rather than an exact selector.  The
composite state stores the finite candidate levels and compatibility edges.
After all requests are known, run backward viability to mark candidates that
can participate in a complete compatible assignment.  If at least one root
candidate is viable, bounded search returns a compatible assignment; if the
output is itself intended to remain relational, return all viable output
candidates.  Soundness follows by first applying the viability theorem to get
the finite compatible family, then applying the soundness derivations of
\(F\) and \(G\) by substitution and cut.  Every loop is bounded by the supplied
finite candidate bounds.
\end{proof}

\subsection{Three levels of output strength}

A theorem should state the weakest output that supports its applications.

\begin{description}[leftmargin=3.5cm,style=nextline]
\item[Relational output.] For every requested finite challenge, return a
      nonempty finite set of correct answers plus compatibility information.
      This needs no global selector.
\item[Approximate functional output.] Return one answer at each precision,
      with a regularity theorem.  This requires a shrinking or determinacy
      modulus as in \cref{sec:shrinking}.
\item[Exact selected output.] Return a distinguished exact representative,
      such as an ordered root, exact sign, normalized eigenvector, or branch.
      This requires a canonical algebraic normalization or a positive
      separation promise.  Without it the request may be obstructed.
\end{description}

The relational form is not a consolation prize.  Many classical theorems are
naturally invariant under permutations, signs, projective scaling, or chart
choices.  Stating the theorem at that invariant level often removes a
constructive obstruction while preserving all mathematically meaningful
content.

\subsection{Challenge trees rather than fixed tables}

An adaptive algorithm does not know its full input demand in advance.  Its
finite demand is a tree: a rational comparison at one stage selects the next
query.  A certificate can include one realized branch, while a uniform
soundness theorem verifies every branch.  The trace state contains the branch
history, and the invariant records that every answer on the history is valid.

If the theorem must be robust under approximate comparisons, branch guards
must overlap.  For example, given an approximation \(r\) to \(x\) within
\(\varepsilon\), one branch may be certified for
\(x\le c+2\varepsilon\) and the other for \(x\ge c-2\varepsilon\).  The
algorithm branches on whether \(r\le c\), but both mathematical branches are
valid in the overlap.  Exact branching is allowed only when an apartness gap
separates \(x\) from \(c\).

\subsection{Extensional interpretation}

A relational presentation represents mathematical content through the family
of accepted finite observations, not through the raw parameter code.  Two
parameters are observationally equivalent when every accepted answer for one
can be matched, at the requested precision, by an accepted answer for the
other and conversely.  This is a finite Hausdorff-style relation at each
precision.  Equality of represented objects is the universal closure over
precisions.

Operations should therefore be proved \emph{extensional}: matched input
observations at the precision requested by a modulus produce matched output
observations.  \Cref{sec:extensionality} gives the generic proof pattern.

\section{Finite witness amalgamation by viability}
\label{sec:amalgamation}

\subsection{Layered finite relations}

Fix a numerical parameter \(p\) and a horizon \(h\).  At level \(k\le h\),
answers are numbers \(a<B(p,k)\).  Two Boolean predicates are given:
\[
 \Good(p,k,a),\qquad
 \Compat(p,k,a,b),
\]
where the second predicate is used only when
\(a<B(p,k)\) and \(b<B(p,k+1)\).  A path of length \(d\) beginning at level
\(k\) is a tuple
\[
 (a_0,\ldots,a_d),
 \qquad a_j<B(p,k+j),
\]
such that every \(a_j\) is good and every adjacent pair is compatible.
Tuples are coded by any fixed primitive-recursive list or bounded-sequence
coding.  The exact coding is immaterial once projection, length, and bounded
folds have their defining equations.

A na\"ive attempt to construct a path chooses a good answer at each level and
hopes that the choices remain extendible.  This can fail even in a finite
binary tree: a locally good node may be a dead end.  Backward viability removes
this failure without using any choice principle.

\subsection{Backward viability}

Define \(\Viable(p,k,d,a)\) recursively in the remaining depth \(d\):
\begin{align*}
 \Viable(p,k,0,a)
   &:= [a<B(p,k)]\wedge \Good(p,k,a),\\
 \Viable(p,k,d+1,a)
   &:= [a<B(p,k)]\wedge \Good(p,k,a)\wedge{}\\
   &\quad \Any_{b<B(p,k+1)}
      \bigl(\Good(p,k+1,b)\wedge
             \Compat(p,k,a,b)\wedge
             \Viable(p,k+1,d,b)\bigr).
\end{align*}
Every connective and bounded existential in this definition is a Boolean fold,
so \(\Viable\) is primitive recursive.  The changing level causes no problem:
the recursion variable is \(d\), while \(k\) is an ordinary parameter.

\begin{theorem}[Finite path--viability equivalence]
\label{thm:path-viability}
For every fixed layered presentation, iPRA proves
\begin{align*}
 \Viable(p,k,d,a)=1
 \quad\Longleftrightarrow\quad&
 \text{there is a coded compatible good path of length }d\\
 &\text{beginning with }a\text{ at level }k.
\end{align*}
The witness path may be bounded by a primitive-recursive code bound depending
only on \(p,k,d\) and the candidate bounds \(B\).
\end{theorem}

\begin{proof}
Induct on \(d\).  At depth zero, the singleton path \([a]\) is valid exactly
when \(a\) is in range and good.  For the successor direction, unfold the
recursion equation for \(\Viable\).  Correctness of the finite \(\Any\)-fold
gives a witness \(b<B(p,k+1)\) satisfying the edge condition and
\(\Viable(p,k+1,d,b)=1\).  The induction hypothesis yields a path beginning
with \(b\); prepend \(a\).  Conversely, a path beginning with \(a,b\) supplies
the required bounded witness \(b\), while its tail supplies viability by the
induction hypothesis.

The induction formula can be Booleanized as equality of two checkers: the
viability bit and the bounded path-existence bit.  Hence the formal proof uses
quantifier-free induction.  The displayed existential statement is recovered
from the search-correctness theorem after the Boolean induction is complete.
\end{proof}

\subsection{A canonical finite path}

Let
\[
 \Next(p,k,d,a)
\defeq \Search_{b<B(p,k+1)}
  \bigl(\Good(p,k+1,b)\wedge\Compat(p,k,a,b)
       \wedge\Viable(p,k+1,d,b)\bigr).
\]
When \(\Viable(p,k,d+1,a)=1\), bounded-search correctness proves that
\(\Next(p,k,d,a)\) is a compatible viable successor.  Primitive recursion on
\(j\le d\), repeatedly applying \(\Next\), gives a canonical path
\(\mathsf{path}(p,k,d,a,j)\).

\begin{corollary}[Least viable path selector]
\label{cor:least-viable-path}
There is a fixed primitive-recursive path constructor such that iPRA proves
\[
 \Viable(p,k,d,a)=1\to
 \mathsf{PathOK}(p,k,d,a,\mathsf{path}(p,k,d,a))=1.
\]
No uniqueness hypothesis is needed.  The least-witness convention is only a
finite implementation choice; the mathematical conclusion remains relational.
\end{corollary}

\begin{proof}
Use quantifier-free induction on the number of generated edges.  The invariant
states that the current node is viable for the number of edges still required
and that the accumulated path is good and compatible.  Search correctness
establishes the successor invariant.  At the final stage, depth-zero viability
supplies goodness of the last node.  All path-prefix assertions are finite
Boolean folds.
\end{proof}

\subsection{From local seriality to finite chains}

Sometimes the input theorem supplies not a precomputed viable root but a
seriality contract:
\begin{align*}
 &\exists a<B(p,0)\;\Good(p,0,a),\tag{I}\label{eq:serial-init}\\
 &\forall k<h\;\forall a<B(p,k)\;
   \bigl(\Good(p,k,a)\to
     \exists b<B(p,k+1)\,[\Good(p,k+1,b)\wedge\Compat(p,k,a,b)]\bigr).
   \tag{E}\label{eq:serial-step}
\end{align*}
For a fixed finite horizon, bounded search converts these assumptions to a
path without countable choice.

\begin{theorem}[Finite dependent choice is primitive recursive]
\label{thm:finite-dependent-choice}
For every horizon \(h\), assumptions \eqref{eq:serial-init} and
\eqref{eq:serial-step} yield a primitive-recursively constructed compatible
good path through levels \(0,\ldots,h\).  The theorem is provable in iPRA.
\end{theorem}

\begin{proof}
Choose the least initial good candidate by bounded search.  At stage \(k<h\),
specialize the seriality premise to the current candidate and choose the least
compatible successor by bounded search.  Induct on \(k\) with the
quantifier-free invariant that the coded prefix has the required length and
its Boolean path checker equals one.  The horizon is a free numerical
parameter; there is no infinite sequence and no appeal to countable choice.
\end{proof}

Backward viability is stronger than this serial construction.  It works when
some good nodes are dead ends, and it supports additional global constraints
that can be tested only after later levels are known.

\subsection{Constraint propagation and witness amalgamation}

Let \(C(p,k,a)\) be an additional Boolean constraint, perhaps a prescribed
residual budget, chart membership, sign bracket, or compatibility with a
second proof module.  Replace \(\Good\) by \(\Good\wedge C\) in the viability
recursion.  This gives the following monotonicity and amalgamation principle.

\begin{theorem}[Finite witness amalgamation]
\label{thm:witness-amalgamation}
Suppose a finite collection of requirements is compiled into one layered
finite relation.  If the initial viability fold is nonempty,
\[
 \Any_{a<B(p,0)}\Viable(p,0,h,a)=1,
\]
then bounded search and the least viable path constructor produce one finite
witness family satisfying all requirements simultaneously.  Adding a new
Boolean requirement can only remove viable nodes; it cannot create a false
path.
\end{theorem}

\begin{proof}
Search supplies a viable initial node.  Apply
\cref{cor:least-viable-path}.  Soundness of conjunction and of every folded
requirement shows that the output path satisfies each component.  For
monotonicity, an induction on remaining depth proves that viability under the
stronger goodness predicate implies viability under the weaker one.  At the
successor step, the same witness edge is used.
\end{proof}

This is a finite constraint-satisfaction theorem.  It applies not only to
linear paths but, after coding the state of a bounded collection of pending
requirements, to finite trees, finite acyclic proof networks, overlapping
chart choices, and simultaneous approximation windows.  The state-space may
be large, but it remains primitive recursive and its correctness proof is the
same Boolean invariant.

\subsection{Full-horizon viability}

The next section needs to know not merely that a node extends forward, but that
it occurs on a complete path from level zero to a specified horizon.  Define
\[
 \FullV(p,h,k,a)=1
\]
to mean that \(k\le h\) and there is a compatible good path
\((a_0,\ldots,a_h)\) with \(a_k=a\).  This is a bounded path-existence
predicate and hence primitive recursive.  It also admits a dynamic-programming
factorization:
\[
 \FullV(p,h,k,a)
 \quad\Longleftrightarrow\quad
 \Reach(p,k,a)\wedge\Viable(p,k,h-k,a),
\]
where \(\Reach\) is the forward Boolean fold saying that \(a\) is reachable
from some good level-zero node.  Both directions are proved by concatenating
or splitting finite paths.

This two-sided notion is crucial.  A merely forward-viable candidate at a high
level need not be connected to the input data.  A merely reachable candidate
may be a dead end.  Full-horizon viability records exactly the candidates that
participate in a complete finite solution.

\subsection{Object-language implementation}

There are two efficient implementations, both suitable for the trace compiler.

\paragraph{Path-code checker.}
Bound a path code, search over all codes, and verify range, goodness,
compatibility, and the distinguished coordinate by finite folds.  This version
has the smallest correctness proof and is useful for certificates.

\paragraph{Dynamic-programming table.}
Build the forward reachable table and backward viable table, each represented
as a bounded bit-vector code.  The local invariant states that the first
\(j\) entries agree with the defining fold.  This version avoids exponential
search and is appropriate for compiled algorithms.

In either implementation, the mathematical theorem is an instance of
\cref{thm:embedded-trace}.  The application module must supply object-language
derivations for Boolean lookup, update, folds, and the local recurrence; the
generic trace theorem then supplies the global soundness derivation.

\section{Shrinking viability and coherent approximate names}
\label{sec:shrinking}

\subsection{The remaining coherence problem}

Finite amalgamation solves every theorem with a fixed finite horizon.  It does
not by itself produce one infinite coherent branch.  Nor should it: an
infinite finitely branching tree can have paths of every finite length while a
uniform choice of an infinite branch may require additional logical or
computational strength, depending on how the tree is presented.

Most mathematical applications need less.  They do not need a stable code for
the same branch at every level; they need observations that become mutually
close.  If every full finite solution has nearly the same observation at the
requested level, then arbitrary finite choices can be determinized without
choosing an infinite branch.  This is the content of the shrinking-viability
theorem.

\subsection{Metric observations}

For each good candidate \(a<B(p,k)\), let
\(\Obs(p,k,a)\) be a rational or coded metric observation.  Let
\(d(x,y)\) be a primitive-recursive rational-valued pseudometric on such
observations.  We require the ordinary finite metric laws as embedded
inequalities.  Exact equality of codes is not required: different candidates
may denote the same approximate object.

Assume a primitive-recursive edge budget \(\delta(p,k)\ge0\) such that
\begin{equation}
 \Good(p,k,a)\wedge\Good(p,k+1,b)\wedge\Compat(p,k,a,b)
 \quad\Longrightarrow\quad
 d(\Obs(p,k,a),\Obs(p,k+1,b))\le\delta(p,k).
 \label{eq:edge-budget}
\end{equation}
Let
\[
 \Tail(p,k,\ell)=\sum_{j=k}^{\ell-1}\delta(p,j)
 \qquad(k\le\ell)
\]
be the finite telescoping budget.

Choose a nondecreasing primitive-recursive horizon function \(H(p,k)\) with
\(H(p,k)\ge k\).  Define the finite viable cloud
\[
 V_{p,k}:=\{a<B(p,k):\FullV(p,H(p,k),k,a)=1\}.
\]
It is a finite, decidable set.  Its canonical representative is
\[
 s(p,k)=\min V_{p,k},
\]
with any fixed default value when the set is empty.  In applications, a
separate hypothesis proves nonemptiness.

\subsection{The shrinking-viability theorem}

\begin{theorem}[Shrinking viability]
\label{thm:shrinking-viability}
Suppose iPRA proves the following numerical contracts for every admissible
parameter \(p\).

\begin{enumerate}[label=(SV\arabic*)]
\item \emph{Nonempty finite solutions:}
      \(V_{p,k}\ne\varnothing\) for every \(k\).
\item \emph{Monotone horizons:}
      if \(k\le\ell\), then \(H(p,k)\le H(p,\ell)\).
\item \emph{Horizontal shrinking:} for all \(a,b\in V_{p,k}\),
      \[
       d(\Obs(p,k,a),\Obs(p,k,b))\le\varepsilon(p,k).
      \]
\item \emph{Edge control:} every compatible edge satisfies
      \eqref{eq:edge-budget}.
\end{enumerate}
Then the primitive-recursive selector
\[
 x(p,k):=\Obs(p,k,s(p,k))
\]
satisfies, for all \(k\le\ell\),
\begin{equation}
 d(x(p,k),x(p,\ell))
 \le \varepsilon(p,k)+\Tail(p,k,\ell).
 \label{eq:shrinking-bound}
\end{equation}
Consequently, if a primitive-recursive modulus \(r(p,n)\) satisfies
\[
 k\ge r(p,n)\to
 \varepsilon(p,k)+\sum_{j=k}^{m-1}\delta(p,j)\le 2^{-n}
 \quad\text{for every }m\ge k,
\]
then \(k\mapsto x(p,r(p,k))\) is a regular Cauchy name with an explicit
primitive-recursive rate.
\end{theorem}

\begin{proof}
Fix \(p\) and \(k\le\ell\).  By (SV1), the selected candidate
\(s(p,\ell)\) lies on a complete good compatible path
\[
 c_0,c_1,\ldots,c_{H(p,\ell)}
 \qquad c_\ell=s(p,\ell).
\]
Here and below the path is obtained by bounded search from the Boolean
\(\FullV\) assertion; no unbounded choice is used.

Let \(b=c_k\).  Since horizons are nondecreasing,
\(H(p,k)\le H(p,\ell)\).  Truncating the displayed path at level \(H(p,k)\)
shows that \(b\in V_{p,k}\).  The selected candidate \(s(p,k)\) also lies in
\(V_{p,k}\), so horizontal shrinking gives
\[
 d(\Obs(p,k,s(p,k)),\Obs(p,k,b))\le\varepsilon(p,k).
\]
Along the segment \(c_k,c_{k+1},\ldots,c_\ell\), edge control and repeated
triangle inequalities give
\[
 d(\Obs(p,k,b),\Obs(p,\ell,s(p,\ell)))
 \le \sum_{j=k}^{\ell-1}\delta(p,j).
\]
One final triangle inequality proves \eqref{eq:shrinking-bound}.

Every operation in this argument is finite.  In the object theory, the path
is the least bounded path certified by \(\FullV\); membership in the two
viable clouds is a Boolean equation; the telescoping inequality is proved by
quantifier-free induction on \(\ell-k\).  Universal introduction over
\(k,\ell\) is performed after that induction.  The final regularity claim is
obtained by substituting the supplied modulus.
\end{proof}

\subsection{Where horizontal shrinking comes from}

The hypothesis that a viable cloud has small diameter is quantitative
uniqueness, not an unexplained compactness principle.  Several standard
sources are reusable.

\paragraph{Isolating intervals or balls.}
If every full finite solution at level \(k\) is certified to lie in the same
rational ball of radius \(\rho_k\), then
\(\varepsilon_k=2\rho_k\).

\paragraph{Strong convexity.}
Suppose candidates satisfy
\(f(a)\le m_k+\eta_k\), while
\(f(x)\ge f(x_*)+\lambda d(x,x_*)^2\).  Then every viable candidate is within
\(\sqrt{\eta_k/\lambda}\) of the unique minimizer, so the cloud diameter is at
most twice that bound.  In a strictly finitary development, the square-root
bound is replaced by a rational dyadic upper bound verified by squaring.

\paragraph{Residual-to-distance estimates.}
If a certified inequality
\[
 d(x,Z)\le C\,R(x)
\]
relates residual to distance from one isolated solution or one small cluster,
and every viable candidate has residual at most \(\eta_k\), then the cloud
diameter is bounded by twice \(C\eta_k\).

\paragraph{Algebraic separation.}
For a polynomial with exact algebraic coefficient data, isolating disks and a
root-separation bound group approximations into disjoint clusters.  Within a
cluster the viable cloud shrinks; at collisions the cluster or root multiset,
not an ordered root, is the stable object.

\paragraph{Contraction certificates.}
If a rational box is mapped into itself by a contraction of factor
\(q<1\), any two candidates surviving \(n\) refinement steps differ by at
most \(q^n\) times the initial diameter plus the verified rounding budget.

These mechanisms are exactly the data absent from a bare qualitative
uniqueness statement.  The theorem refuses to infer a modulus from uniqueness
alone.

\subsection{Nested balls as a special case}

\begin{corollary}[Nested-ball name]
\label{cor:nested-ball-name}
Suppose a certificate at level \(k\) is a rational center \(c_k\) and radius
\(r_k\), every full path chooses nested balls, and every viable ball at level
\(k\) intersects one common rational ball of radius \(r_k\).  If
\(r_k\le 2^{-k-3}\) and compatible centers move by at most
\(2^{-k-3}\), then the least full-viable center, shifted by a fixed number of
precision levels, is a regular name.
\end{corollary}

\begin{proof}
Intersecting balls of radius \(r_k\) have centers at distance at most
\(2r_k\), so take \(\varepsilon_k=2r_k\).  The edge budgets form a geometric
series.  Apply \cref{thm:shrinking-viability} and absorb the constant factor by
shifting the requested precision.
\end{proof}

\subsection{Exact choice under a positive gap}

Shrinking gives approximate coherence.  Exact labels can be recovered only
when the labels are separated.

\begin{corollary}[Eventual exact label stability]
\label{cor:exact-label-gap}
Suppose candidates carry labels in a finite set, candidates with different
labels have observational distance at least \(g>0\), and the right side of
\eqref{eq:shrinking-bound} is eventually less than \(g\).  Then the selected
labels are eventually constant.  A primitive-recursive search for the first
precision where the certified error is below \(g\) yields an exact label.
\end{corollary}

\begin{proof}
Two selected candidates at sufficiently high precisions are closer than
\(g\); therefore they cannot have different labels.  The threshold is found
from the supplied numerical modulus, not from a qualitative assertion that a
gap exists.
\end{proof}

This is the correct interface for root ordering, pivot choice, chart choice,
sign, and projective orientation: approximate or relational data are
unconditional; exact labels require a positive rational separation witness.

\subsection{What happens without shrinking}

If (SV3) is unavailable, the framework still yields strong conclusions.
For every requested finite collection of precisions, finite amalgamation
produces a compatible family.  Downstream theorems that are stable under all
such families can consume the relation directly.  Examples include:

\begin{itemize}
\item existence of a root cluster meeting a residual bound;
\item the minimum value of a function without a selected minimizer;
\item approximate rotation data without a convergent axis;
\item a finite cover or parity witness in topology;
\item all representatives passing a finite test suite.
\end{itemize}

The absence of a coherent selector is then a property of the statement, not a
failure of the encoding.  This distinction prevents hidden uses of choice,
LLPO, or discontinuous branch selection.

\section{Guarded extensionality and precision algebra}
\label{sec:extensionality}

\subsection{Observation equivalence rather than quotient types}

The object language of iPRA is first order over natural numbers.  It need not
contain quotient types for real numbers, probability measures, projective
points, or functions.  Instead, a mathematical object is handled through a
fixed numerical presentation and an observational equivalence relation.

For a functional presentation with evaluator \(a(p,k)\), define
\[
 p\equiv_X q
 \quad:\Longleftrightarrow\quad
 \forall k\;\Close_X(p,q,k),
\]
where \(\Close_X\) is a primitive-recursive finite inequality.  For a
relational presentation, the natural notion is Hausdorff-style finite
closeness of viable answer clouds, or equivalently agreement of all stable
finite tests.  In either case, the global equivalence is not inserted into an
induction formula.  At a fixed output challenge, only finitely many instances
of \(\Close_X\) are required.

\begin{principle}[Extensional proof discipline]
To prove that a numerical construction \(F\) is an operation on represented
objects, first choose a primitive-recursive precision transformer
\(\mu_F\).  Prove the fixed-precision implication
\[
 \bigwedge_{q\in Q_F(p,k)}\Close_X(p,p',q)
 \longrightarrow
 \Close_Y(F(p),F(p'),k),
\]
where \(Q_F(p,k)\) is the finite query table generated by \(\mu_F\) and the
bounded trace.  Only afterward introduce the universal quantifier over
\(k\).
\end{principle}

This makes extensionality a theorem with a visible numerical budget rather
than a hidden appeal to continuity.

\subsection{Composition of precision contracts}

A certified operation \(F:X\to Y\) has three numerical interfaces:

\begin{enumerate}
\item an output algorithm or output relation;
\item a demand function \(\mu_F(p,k)\), possibly returning a finite table of
      input challenges; and
\item an error transformer \(E_F\) proving that input errors within the
      demanded budget yield the requested output error.
\end{enumerate}

If \(G:Y\to Z\), the composite demand is obtained by first evaluating the
finite metadata needed for \(G\), then substituting every challenge requested
by \(G\) into \(F\)'s demand.  For scalar precision transformers this is the
familiar formula
\[
 \mu_{G\circ F}(p,k)=\mu_F(p,\mu_G(F_0(p),k)).
\]
For adaptive query tables, concatenate the proof-producing traces.  Error
budgets compose by ordinary rational inequalities.

\begin{theorem}[Precision-contract composition]
\label{thm:precision-composition}
Certified extensional transducers are closed under sequential composition,
finite products, and bounded case distinction.  The composite precision map,
trace bound, guard, and proof term are primitive-recursively computable from
the component data.
\end{theorem}

\begin{proof}
Sequential composition is the phase construction of the trace compiler.  A
finite product pairs states, unions the finite query tables, and takes the
maximum demanded precision.  A bounded case split computes every guard as a
Boolean and uses a phase tag to select the corresponding invariant.  In each
case the final extensionality implication is obtained by cut from the
component implications.  Construction of the resulting \(\Prf\) term is
structural recursion on the finite transducer expression.
\end{proof}

\subsection{Totalization and guards}

Primitive-recursive function symbols are total, whereas familiar
mathematical operations are often partial.  The correct solution is not to
pretend that the denominator, pivot, chart, logarithm argument, or derivative
is always valid.  Define an arbitrary totalized numerical operation and state
its correctness under a Boolean guard.

For example, a coded rational inverse may return zero when the numerator is
zero.  Its theorem has the form
\[
 \mathsf{Apart}(p,d)=1\to
 \mathsf{InvContract}(p,d,k,\mathsf{invApprox}(p,d,k))=1.
\]
The output outside the guard has no mathematical interpretation.  Because the
operation itself is total, it is a legitimate \(\PR\) symbol; because the
soundness theorem exposes the guard, no partiality is concealed.

\begin{principle}[No unguarded partial operation]
Every use of division, normalization, pivoting, branch inversion, logarithm,
argument, local chart, or reciprocal must consume a finite positive margin.
The margin may be an input certificate or may be produced by an earlier
verified module.  A mere negated equality is not silently converted to a
positive rational lower bound.
\end{principle}

\subsection{Overlapping branches}

Exact branching on an approximate comparison is unstable.  A robust
alternative uses overlap.  Suppose two formulas \(F_0(x)\) and \(F_1(x)\) are
both valid on an overlap region
\([	au-\eta,\tau+\eta]\), while the first is valid to the left and the second
to the right.  Approximate \(x\) within \(\eta/2\):

\begin{itemize}
\item if the approximation is below \(\tau\), return branch zero;
\item otherwise return branch one.
\end{itemize}

If the comparison is wrong because \(x\) lies near \(\tau\), both branches are
valid.  The construction is total and continuous at the level of its
contract, even though the returned branch bit may fluctuate.

\begin{theorem}[Robust overlapping branch]
\label{thm:overlapping-branch}
Given finite proofs of left validity, right validity, overlap validity, and an
approximation error below half the overlap width, iPRA proves that the branch
selected by the rational comparison satisfies the common output contract.
\end{theorem}

\begin{proof}
There are two decidable rational cases.  In the selected-left case, either the
true value is left of the overlap endpoint, where left validity applies, or
the approximation error places it in the overlap.  The selected-right case is
symmetric.  The proof is finite rational arithmetic and can be Booleanized.
\end{proof}

\subsection{Exact branching under a gap}

If an input certificate gives
\[
 x\le \tau-g\quad\text{or}\quad x\ge \tau+g
\]
for a positive rational \(g\), an approximation within \(g/3\) determines the
correct side exactly.  Thus exact sign, orientation, chart, pivot, and root
labels are available under quantitative separation.

\begin{corollary}[Exact branch from a positive gap]
\label{cor:exact-branch-gap}
A supplied disjunctive gap certificate and a sufficiently accurate rational
observation yield a primitive-recursive exact branch selector whose
correctness is provable in iPRA.
\end{corollary}

The disjunction is important.  From \(x\ne0\) alone, intuitionistic logic does
not provide either a sign or a rational lower bound for \(|x|\).

\subsection{Congruence for relational outputs}

A relational operation is extensional when close input clouds produce close
output clouds in both directions: every valid output from one input has a
compatible nearby output for the other, and conversely.  This is a finite
back-and-forth condition at each challenge.  It composes by witness
amalgamation.

\begin{theorem}[Relational congruence]
\label{thm:relational-congruence}
Let \(F:X\rightsquigarrow Y\) be a certified relational transducer.  Suppose
input parameters \(p,p'\) satisfy the finite closeness table requested for an
output challenge \(k\).  Then every accepted output answer for \(p\) has a
\(Y\)-compatible accepted output answer for \(p'\), and conversely, provided
the transducer's local back-and-forth obligations are derivable.
\end{theorem}

\begin{proof}
Run the two finite traces in product.  Add the input closeness facts to the
Boolean invariant and pair each nondeterministic answer with a finite set of
compatible answers on the other side.  Backward viability removes dead ends.
The final product invariant gives the two output witnesses.  All searches are
bounded, so the argument is an instance of finite witness amalgamation and
embedded trace reflection.
\end{proof}

This theorem is the appropriate substitute for quotient congruence when
objects have non-canonical representatives.  A deterministic operation is the
special case in which each output cloud is a singleton.

\subsection{Normalization is a theorem, not an identity assumption}

Codes for rationals, polynomials, projective points, matrices, and finite sets
often admit multiple representations.  The framework uses one of two safe
patterns:

\begin{enumerate}
\item choose a primitive-recursive canonical normalization and prove that it
      preserves observations; or
\item retain non-canonical codes and prove every operation respects the
      observational relation.
\end{enumerate}

The first pattern is convenient for signed integers and positive-denominator
rationals.  The second is often better for projective or clustered data, where
canonical choice would require unstable comparisons.  In neither pattern is
code equality confused with mathematical equality.

\section{A quantitative library for analysis and approximation}
\label{sec:analysis-library}

This section records the shared analytic lemmas that should be formalized once
and reused.  The purpose is not to recreate a set-theoretic real-analysis
library inside iPRA.  It is to isolate finite rational contracts sufficient for
the theorem packets.

\subsection{Dyadic and rational names}

A fast rational name is a fixed evaluator \(a(p,k)\in\Q\) satisfying
\begin{equation}
 |a(p,k)-a(p,\ell)|\le 2^{-k}+2^{-\ell}.
 \label{eq:regular-name}
\end{equation}
Rationals are coded by signed numerators and positive denominators, or by
normalized dyadics where convenient.  Absolute value, cross-multiplied order,
addition, multiplication, and finite sums are primitive recursive on the
codes.  Denominator positivity is part of every rational-domain guard.

The relation \eqref{eq:regular-name} is a first-order schema for each fixed
evaluator symbol.  It does not quantify over functions.  Equality of two
names is expressed by
\[
 |a(p,k)-b(q,\ell)|\le 2^{-k}+2^{-\ell}
\]
for all \(k,\ell\), or by an equivalent one-precision formulation with a
constant shift.

\subsection{Addition, multiplication, and inverse}

Addition needs only a constant precision shift.  Define
\[
 c(p,q,k)=a(p,k+2)+b(q,k+2).
\]
Two applications of \eqref{eq:regular-name} and the triangle inequality show
that \(c\) is again regular.  Negation is immediate.

Multiplication additionally needs a magnitude contract.  If all queried
samples of \(a\) and \(b\) have absolute value at most \(M\), then
\[
 |a_kb_k-a_\ell b_\ell|
 \le M|a_k-a_\ell|+M|b_k-b_\ell|.
\]
A primitive-recursive shift \(r(M,k)\) making the right side at most the
requested dyadic error yields a multiplication name.  The magnitude bound may
be supplied globally or computed from a coarse sample plus the regularity
radius.

For inverse, suppose a positive rational \(d\) certifies
\(|a_k|\ge d\) at every precision queried by the algorithm.  Then
\[
 \left|\frac1{a_k}-\frac1{a_\ell}\right|
 =\frac{|a_k-a_\ell|}{|a_ka_\ell|}
 \le d^{-2}|a_k-a_\ell|.
\]
Choose the input precision so that the final expression meets the output
budget.  Numerically, inverse is totalized; correctness is asserted only under
the apartness guard.  Division is multiplication followed by guarded inverse.

\begin{theorem}[Arithmetic of guarded regular names]
\label{thm:regular-name-arithmetic}
For every fixed family of input evaluators, iPRA proves regularity and
extensionality of addition, negation, multiplication under a magnitude bound,
and inverse under a positive apartness bound.  The required precision
transformers are primitive recursive and compose.
\end{theorem}

\begin{proof}
Fix an output precision and instantiate the finitely many input regularity,
magnitude, and apartness facts used at that precision.  Clear positive
rational denominators and prove the resulting integer inequalities.  Each
calculation is a bounded trace whose invariant is the current rational error
budget.  Apply the trace compiler, then universally quantify the output
precision.  The displayed elementary inequalities give the local soundness
proofs.
\end{proof}

\subsection{Finite sums, products, and polynomial evaluation}

A finite sum carries the invariant
\[
 \left|\sum_{i<j}\widetilde x_i-
       \sum_{i<j}x_i\right|
 \le \sum_{i<j}e_i.
\]
A finite product carries a magnitude bound and the telescoping identity
\[
 \prod_{i<n}x_i-\prod_{i<n}y_i
 =\sum_{j<n}(x_j-y_j)
   \left(\prod_{i<j}x_i\right)
   \left(\prod_{j<i<n}y_i\right).
\]
Both are proved by quantifier-free induction on the finite length after all
error quantities have been rationally encoded.  Horner evaluation gives a
particularly compact polynomial trace and error recurrence.  Matrix
multiplication, determinants by a permutation sum or fraction-free
elimination, truncated power series, and finite convolutions are instances of
the same library.

\subsection{Finite-net transfer}

Let \(K(p)\) be a represented compact domain only in the following concrete
sense: for each rational radius \(\rho\), a primitive-recursive function
returns a finite list \(N(p,\rho)\), and a certificate proves
\[
 \forall x\in K(p)\;\exists y\in N(p,\rho)\; d(x,y)\le\rho.
\]
The quantifiers over represented points are interpreted through the finite
presentation used by the theorem; in a fully embedded statement, coverage is
a certificate-relative first-order formula.

Let \(f\) have a modulus \(\omega\): if \(d(x,y)\le\rho\), then
\(|f(x)-f(y)|\le\omega(\rho)\).

\begin{theorem}[Finite-net transfer]
\label{thm:finite-net-transfer}
If \(N\) is a \(\rho\)-net, \(f(y)\ge L\) for every \(y\in N\), and
\(\omega(\rho)\le\eta\), then \(f(x)\ge L-\eta\) throughout the represented
domain.  Likewise, a grid point whose value is within \(\eta\) of the finite
grid minimum is within \(2\eta\) of the global infimum.
\end{theorem}

\begin{proof}
For an arbitrary represented \(x\), the finite coverage certificate and
bounded search give a net point \(y\) with \(d(x,y)\le\rho\).  The modulus
gives \(f(x)\ge f(y)-\eta\ge L-\eta\).  For minimization, compare the selected
grid point to a net point near an arbitrary \(x\), then take the finite minimum
using bounded search.  No abstract compactness or completed infimum object is
used.
\end{proof}

The minimum \emph{value} needs no unique minimizer.  A coherent minimizer name
requires a quantitative uniqueness or shrinking-viability hypothesis.

\subsection{Riemann sums and integration}

An integral theorem is represented by a quadrature algorithm plus two finite
contracts:

\begin{enumerate}
\item a modulus controlling the change of the integrand within each cell; and
\item a tail or truncation estimate if the interval or series is unbounded.
\end{enumerate}

For a bounded interval and a mesh of width \(\rho\), upper and lower rational
sums differ by at most interval length times \(\omega(\rho)\).  The trace
computes the finite sum, while the invariant accumulates rounding and mesh
errors.  Additivity, linearity, and change of variables for explicitly
monotone rational maps reduce to finite sum rearrangements and modulus bounds.

The fundamental theorem of calculus should therefore be split into finite
statements.  For the derivative-to-integral direction, telescope finite
increments and control each local remainder.  For the integral-to-derivative
direction, compare a difference quotient to cell averages using a continuity
modulus.  A theorem packet that states merely that a derivative is continuous
has not yet supplied the numerical input required by the iPRA proof.

\subsection{Series and analytic functions}

A series-defined function is a fixed finite-sum evaluator together with a
primitive-recursive tail bound
\[
 \left|\sum_{j\ge N}t_j(p,z)\right|\le T(p,z,N).
\]
All algebraic identities are first proved for finite truncations.  Tail bounds
then transfer them to the represented function at a fixed requested
precision.  This pattern covers exponential, sine, cosine, logarithm on a
guarded domain, Euler products over bounded primes, and many special-function
identities.

Local analytic inversion uses a finite certificate: a center, rational box,
derivative or Jacobian gap, contraction constant, coefficient or derivative
bounds, and a radius.  Finite iteration constructs approximations; the
contraction estimate gives horizontal shrinking.  Singular branch discovery
is a separate producer problem and must not be hidden in the local theorem.

\subsection{Approximate extrema}

For a finite rational grid, bounded search returns a least residual or least
value point.  Let \(m_k\) denote the finite minimum.  The finite-net theorem
turns \(m_k\) into a name for the infimum value.  To obtain a point:

\begin{itemize}
\item retain the relation of all \(\eta_k\)-near-minimizers; or
\item provide a quantitative uniqueness estimate, such as strong convexity,
      and apply shrinking viability.
\end{itemize}

This resolves a common constructive mismatch.  The operation
``take the infimum value'' is often stable, while ``take the unique point where
it is attained'' needs more data.

\subsection{Finite probability laws}

A finite rational law is a list of outcomes and nonnegative rational weights
summing to one.  Expectation, variance, characteristic functions at rational
arguments, convolution, and distribution values at rational thresholds are
finite sums.  Approximation of an unbounded or continuous law requires
explicit tail and quantization moduli.  Approximation of a discontinuous CDF
at a threshold additionally requires an anti-concentration or boundary-mass
bound.

Thus a constructive central-limit statement has a transparent shape:
finite rational laws, a finite convolution or characteristic-function trace,
a quantitative error estimate, and a rational threshold protected by a
boundary budget.  Higher-type measures are unnecessary unless they are given
through a fixed finite-observation evaluator with the corresponding moduli.

\subsection{Precision ledgers}

Every analytic module should export a machine-readable precision ledger:

\begin{longtable}{@{}p{0.24\textwidth}p{0.68\textwidth}@{}}
\toprule
Field & Meaning\\
\midrule
\endhead
Output challenge & Requested error, region, derivative order, or test.\\
Input queries & Exact finite list or adaptive trace of samples.\\
Magnitude data & Bounds needed for products, tails, or matrix operations.\\
Apartness data & Positive margins needed by inverses and exact branches.\\
Truncation bound & Number of series terms, grid cells, or iterations.\\
Roundoff budget & Rational allocation to every finite operation.\\
Analytic remainder & Tail, modulus, contraction, or residual estimate.\\
Final inequality & Arithmetic implication proving the output contract.\\
\bottomrule
\end{longtable}

The ledger is not auxiliary documentation.  It is the numerical content from
which the fixed-precision object-language proof is built.

\section{Domain libraries built on the common kernel}
\label{sec:domain-libraries}

The framework is useful only if domain proofs share substantial infrastructure.
This section gives the intended library boundaries and the preferred statement
forms for the main areas represented in the repository.

\subsection{Exact finite algebra}

Polynomial and matrix identities should be proved first over coded coefficient
rings by finite sums and products.  A polynomial is a bounded coefficient list;
addition, multiplication, derivative, composition, and evaluation are
primitive-recursive list operations.  Equality is coefficientwise equality,
checked by a finite fold.

The generic matrix layer contains dimensions, entry lookup, matrix products,
permutation or fraction-free determinants, minors, and polynomial matrices.
Cayley--Hamilton is best stated as a uniform finite identity for the generic
matrix, then specialized by coefficient substitution.  This avoids analytic
or field assumptions.  Trace identities, adjugate formulas, and finite linear
recurrences can use the same coefficientwise kernel.

Factorizations and radical expressions are certificate-relative unless a
producer has been formalized.  A factorization certificate consists of a
finite factor list, degree and leading-coefficient data, and a coefficientwise
product identity.  Its verifier is small and uniform.  Irreducibility and
Galois-theoretic impossibility results need finite permutation/group
certificates plus a fully specified arithmetization; a semantic field-theory
proof in ambient Lean is not yet an iPRA derivation.

\subsection{Roots, branches, and algebraic functions}

For approximate roots, the stable primary output is one of:

\begin{enumerate}
\item a finite multiset of isolating boxes with multiplicities;
\item a root cluster whose residual and diameter meet a requested budget;
\item a relational family of approximate roots; or
\item an individual root name under a separation or simple-root certificate.
\end{enumerate}

Ordering roots is deferred until a positive separation witness is available.
Local implicit or inverse functions consume a derivative gap, a rational box,
and a contraction or majorant certificate.  Puiseux expansions from exact
algebraic coefficient codes may use finite Newton-polygon and splitting
certificates.  The stronger demand for exact Puiseux coefficients from
arbitrary Cauchy coefficient names is not a library goal; the repository
contains a computability obstruction even for a fixed evaluator family.

\subsection{Number theory and combinatorics}

Divisibility, primality below a bound, valuations, finite divisor sums,
M\"obius inversion, totients, von Mangoldt weights, binomial coefficients,
Bell recurrences, and congruences are all exact finite arithmetic.  Their
preferred proofs use Boolean divisibility predicates and double-counting
bijections between bounded coded sets.

Asymptotic statements are expressed by explicit inequalities with a free
numerical parameter and a displayed threshold function.  A proof of
\(f(n)\sim g(n)\) becomes
\[
 \forall r\;\forall n\ge N(r)\quad
 |f(n)-g(n)|\le 2^{-r}|g(n)|,
\]
with a suitable integer formulation after clearing denominators.  The
threshold \(N\) must be primitive recursive or supplied as a certificate.
Divergence is stated by a primitive-recursive witness function: for every
height \(K\), a displayed cutoff makes the finite partial sum exceed \(K\).

Analytic number theory should separate a uniform finite verifier from the
producer of difficult constants.  Zero-free regions, exceptional-zero data,
contour bounds, and explicit estimates can be represented by rational tables
whose checker verifies all local inequalities.  This does not solve the
problem of producing effective constants in results where current proofs are
ineffective; it makes the residual obligation exact.

\subsection{Diophantine geometry}

A rational point is encoded by a normalized integer tuple with a height.
Polynomial satisfaction is decidable after clearing denominators.  Once a
height bound \(H\) is supplied, complete enumeration is a bounded search.
Conversely, from a complete finite list one obtains a height bound by taking a
finite maximum.

\begin{theorem}[Height bound--complete list equivalence]
\label{thm:height-list}
For every fixed primitive-recursive point predicate \(P(p,x)\) and height
function \(h(x)\), iPRA proves:

\begin{enumerate}
\item a bound \(H\) satisfying \(P(p,x)\to h(x)\le H\) yields a complete
      finite list of all points by bounded enumeration; and
\item a complete finite list yields the height bound equal to the maximum
      height of its members.
\end{enumerate}
\end{theorem}

\begin{proof}
For the first direction, enumerate every normalized code of height at most
\(H\), retain exactly those passing the Boolean predicate, and prove
completeness by specializing the bound hypothesis.  For the second, compute
the finite maximum.  If \(P(p,x)\), list completeness gives an index at which
\(x\) occurs, so its height is below the maximum.  Both directions are finite
folds and bounded searches.
\end{proof}

This theorem sharpens the boundary in the Faltings, Siegel, and Mordell--Weil
packets.  A verifier and enumerator can be completed independently of the
deep theorem that produces the height bound, finite cover, or rational
covering representatives.  A cardinality bound alone supplies neither.

Descent arguments should be packaged as a well-founded numerical measure,
a primitive-recursive descent step, and a proof that every nonterminal state
strictly decreases.  Primitive-recursive descent then yields the terminal
certificate.  Covering descents are relational when several next coverings
are possible; backward viability handles finite compatibility constraints.

\subsection{Geometry}

Coordinate geometry is naturally finite.  Points, lines, circles, conics, and
projective incidences are coded by rational or represented coordinates.
Determinantal incidence identities, Ceva--Menelaus-type products, power of a
point, and many triangle formulas reduce to polynomial identities plus
nondegeneracy guards.

Angle-based theorems require more care.  Prefer dot products, determinants,
complex phases, quaternions, or sine/cosine pairs to a globally chosen angle.
Normalization needs a positive norm gap.  For rotations, exact quaternion data
can remain meaningful when the vector part vanishes, whereas a unit axis at
the identity is discontinuous.  Approximate axes may be returned relationally;
exact axes require separation from the identity or a supplied nonzero vector.

Spherical and trigonometric identities should be split into exact polynomial
identities in sine/cosine variables and analytic existence/branch statements.
The first belong to the algebraic checker.  The second consume quadrant,
nonvanishing, or interval certificates.  This avoids importing classical
inverse-function conventions into the algebraic part.

\subsection{Topology}

Abstract compactness is replaced by finite combinatorics.  A quantitative
Jordan or invariance-of-dimension theorem should specify:

\begin{itemize}
\item a rational grid scale obtained from a continuity and separation modulus;
\item a finite cell or cubical complex;
\item a Boolean labeling rule;
\item a local parity or Sperner invariant; and
\item a finite path, chain, or parity certificate producing the conclusion.
\end{itemize}

The global proof is quantifier-free induction over cells or a finite trace.
Refinement from one grid to the next is a relational presentation; a coherent
geometric point is extracted only when nested cells shrink quantitatively.

\subsection{Analysis and special functions}

Calculus, Gamma and zeta identities, residues, and asymptotics should be
organized around finite sums, rational contours or meshes, and explicit
remainders.  A complex contour is a finite polygonal path or a fixed evaluator
with a modulus.  Cauchy and residue statements become finite telescoping
identities plus bounds on local approximation errors.  Winding and zero counts
are exact only under a boundary-apartness certificate.

Special functions are best given by a fixed family of finite approximants and
one or more interchangeable certificates: series tails, product tails,
integral quadrature errors, or functional-equation residuals.  Identities are
proved at finite precision through whichever representation has the easiest
error ledger.  Extensionality then shows that the alternatives describe the
same represented value.

\subsection{Probability}

The base probability library consists of finite rational laws, products,
convolutions, moments, and rational test functions.  Limit theorems are
quantitative approximation theorems.  Tail truncation, lattice spacing,
variance lower bounds, third-moment bounds, and anti-concentration are explicit
inputs or derived certificates.  Multivariate laws use finite rational vectors
and finite nets of test directions.  Martingale statements use bounded finite
filtrations and transition tables.

This formulation captures the effective finite content of central-limit
arguments while avoiding an unsupported higher-type measure object.  A more
general represented measure can be added later as an FRP whose challenges are
finite rational test functions and whose answers are certified approximate
integrals.

\subsection{A uniform package boundary}

Every domain theorem should be split into four files or namespaces:

\begin{enumerate}
\item \emph{Data}: codes, guards, evaluators, candidate bounds, and challenges;
\item \emph{Finite kernel}: Boolean folds, arithmetic identities, local
      invariants, and certificate checking;
\item \emph{Producer}: optional algorithm constructing the certificate or
      viable relation; and
\item \emph{Export}: the exact \(\Prf\,\PRAAx\) theorem and its transport by
      \code{pra\_iff\_native}.
\end{enumerate}

A packet may honestly complete the finite kernel while leaving the producer
open.  This is materially stronger and more reusable than an all-or-nothing
outline, provided the status is stated explicitly.

\section{Resolution matrix for the reported difficulties}
\label{sec:resolutions}

The following matrix distinguishes a general resolution from a conditional
interface and from a genuine retained barrier.  ``Resolved'' means that the
logical or representational obstacle has a reusable finitary theorem; it does
not mean that every theorem-specific estimate has been proved.

\begin{longtable}{@{}p{0.15\textwidth}p{0.24\textwidth}p{0.35\textwidth}p{0.18\textwidth}@{}}
\toprule
Area & Reported difficulty & Framework treatment & Status\\
\midrule
\endhead

Object theory & An ambient Lean theorem or passing executable check is treated
as an iPRA proof. & Require local premises and the endpoint to have type
\(\Gamma\vdash[\PRAAx]\varphi\).  The proof-producing trace compiler builds
this derivation; \code{pra\_iff\_native} is invoked only afterward. & Generic
human construction complete; Lean implementation pending.\\
\addlinespace

Induction & A completed all-precision real assertion is not quantifier-free. &
Fix the requested precision, instantiate a finite query table, induct on a
Boolean trace invariant, then universally introduce precision. & Resolved as a
general schema.\\
\addlinespace

Bounded quantifiers & Finite assertions are written with logical quantifiers
inside an invariant. & Compile bounded universal and existential quantifiers
to \(\All\), \(\Any\), and least-witness folds; prove fold correctness once. &
Resolved.\\
\addlinespace

Higher-type input & A universal total evaluator for arbitrary primitive-
recursive program codes is assumed. & Use fixed evaluator schemas; otherwise
supply a clocked finite trace and verify it. & Stronger universal interface is
impossible by diagonalization.\\
\addlinespace

Approximate witnesses & Separate witnesses exist at every precision but no
coherent choice is justified. & Use an FRP, backward viability, and finite
witness amalgamation.  Apply shrinking viability only with a quantitative
cloud-diameter modulus. & Finite case resolved; infinite name conditional on
quantitative determinacy.\\
\addlinespace

Root ordering & Roots collide, so individual labels are unstable. & Return a
multiset, clusters, or relational approximations.  Recover labels only under a
root-separation certificate. & Resolved formulation.\\
\addlinespace

Puiseux branches & Exact branch coefficients are requested from arbitrary
Cauchy-presented coefficients. & Use exact algebraic coefficient codes or
explicit splitting/nondegeneracy certificates; otherwise retain approximate
or multivalued data. & Unrestricted producer remains obstructed by a halting
reduction.\\
\addlinespace

Rotation axis & A unit axis is requested uniformly, including the identity. &
Return quaternion data or approximate axes; exact unit axis only with a gap
from the identity or a supplied nonzero axis. & Unrestricted exact selector
remains LLPO-obstructed.\\
\addlinespace

Partial operations & Division, normalization, pivots, and charts are used from
mere nonzeroness. & Totalize numerical operations and consume an explicit
positive rational guard.  Use overlapping branches near boundaries. &
Resolved under guards; nonzero-to-apartness conversion not assumed.\\
\addlinespace

Compactness & A finite subcover, minimizer, or convergent subsequence is chosen
abstractly. & Supply primitive-recursive finite nets, coverage certificates,
moduli, and bounded minimization.  Use shrinking viability for coherent points.
& Resolved for effective presentations.\\
\addlinespace

Metastability & Stable finite windows are mistaken for a global Cauchy rate. &
Keep the finite-window theorem as its own output.  A Cauchy name requires an
additional modulus or shrinking condition. & Distinction resolved; general
conversion is false.\\
\addlinespace

Faltings/Siegel & Finiteness or a cardinality bound is used as a search cutoff.
& Separate a uniform finite verifier/enumerator from the height-bound producer;
use the height-bound--complete-list equivalence. & Verifier resolved; general
height producer remains open.\\
\addlinespace

Mordell--Weil & Finite Selmer data are treated as rational covering
representatives. & Make representatives and local compatibility a checked
certificate; use relational descent and viability. & Uniform representative
producer remains open.\\
\addlinespace

Prime distribution & Classical proofs contain ineffective constants or
exceptional-zero alternatives. & State a finite analytic-certificate verifier
with correct quantifier order; prove the arithmetic implication from accepted
constants. & Verifier architecture resolved; full constant producer remains
theorem-specific and sometimes ineffective.\\
\addlinespace

Probability & General measures and discontinuous CDF tests obscure finite
content. & Start with finite rational laws; add tail, quantization, variance,
and anti-concentration moduli; use rational thresholds. & Resolved for
effective presentations.\\
\addlinespace

Topology & Continuous compact spaces are used without numerical control. &
Compile to finite rational grids, cubical complexes, parity/Sperner
certificates, and refinement relations. & Representation resolved;
theorem-specific combinatorics remain.\\
\addlinespace

Analytic branches & A local inverse or zero is selected without a derivative
or boundary gap. & Consume a contraction, Jacobian, Rouch\'e, or boundary
apartness certificate.  Return clusters when the gap vanishes. & Resolved
conditionally; singular chart production separate.\\
\addlinespace

Extensionality & Equality of codes is confused with equality of represented
objects. & Use observational setoids or finite back-and-forth relations; prove
operations congruent with explicit precision transformers. & Resolved.\\
\bottomrule
\end{longtable}

\subsection{Three kinds of residual work}

After applying the framework, every remaining gap should fall into one of three
categories.

\paragraph{Finite-kernel engineering.}
The required theorem is elementary and finitary, but the object-language
library for tuples, rationals, finite folds, or error arithmetic has not yet
been implemented.  This is substantial formalization work, but its statement
and induction invariant are known.

\paragraph{Quantitative mathematical estimate.}
The representation is adequate, but a theorem-specific modulus, separation
bound, zero-free estimate, contraction constant, or combinatorial invariant is
missing.  This is the genuine mathematics of the application.

\paragraph{Producer problem.}
A small checker is complete, but no uniform primitive-recursive procedure is
known for producing the accepted certificate.  Height bounds, rational
covering representatives, effective analytic constants, and singular branch
data often live here.  The packet should expose this producer as a separate
theorem rather than weakening the verifier's status.

This classification is itself useful: it prevents a missing library lemma
from being described as a foundational obstruction, and prevents a deep
ineffective theorem from being described as routine encoding.

\section{Lean formalization architecture}
\label{sec:lean-plan}

\subsection{Target environment}

The formalization should live inside the existing ProveIt Lean workspace and
import the checked modules rather than copying a second kernel:

\begin{lstlisting}[style=lean]
import ConciseIPRA.PR
import ConciseIPRA.NativeBiInterpretation

namespace LeanProofs.ConciseIPRA.Finitary
open Syn
open scoped Fm
open PTm (app zero succ)
\end{lstlisting}

The source-accurate interfaces used by the plan are:

\begin{itemize}
\item \code{PR}, \code{PTm}, and \code{PRAAx};
\item \code{PRA.precZero}, \code{PRA.precSucc}, and \code{PRA.qfInd};
\item the constructors and structural lemmas of \code{Prf}, including
      substitution, weakening, cut, universal introduction/elimination, and
      equality substitution; and
\item \code{pra\_iff\_native} for final transport.
\end{itemize}

The development should not start with Mathlib real numbers.  Ambient finite
models are useful for testing, but the primary endpoint of each generic lemma
is a value of type \code{Prf PRAAx Gamma phi}.

\subsection{Canonical code library}

The first module, proposed as \code{Finitary/Codes.lean}, chooses exactly one
code for each recurring data type:

\begin{enumerate}
\item Booleans normalized to zero and one;
\item pairs and bounded tuples;
\item bounded lists and bit vectors;
\item signed integers;
\item positive-denominator rationals and dyadics;
\item finite maps represented by bounded tables.
\end{enumerate}

For every operation there are three layers:

\begin{enumerate}
\item a \code{PR} symbol;
\item a \code{PTm} notation/helper for applying the symbol; and
\item embedded defining/correctness derivations.
\end{enumerate}

Parallel encodings should be avoided.  Most proof-packet integration failures
would be amplified if each packet introduced its own rational or list code.

\subsection{Boolean folds and bounded search}

The next modules are \code{Boolean.lean}, \code{Fold.lean}, and
\code{Search.lean}.  Suggested theorem signatures are schematic generators:

\begin{lstlisting}[style=lean]
/-- Object-language proof that a fixed PR Boolean fold is correct. -/
def allFoldCorrect (b : PR (r + 1)) :
  [] |-[PRAAx] allFoldSpec b := ...

def anyFoldCorrect (b : PR (r + 1)) :
  [] |-[PRAAx] anyFoldSpec b := ...

def leastSearchCorrect (b : PR (r + 1)) :
  [] |-[PRAAx] leastSearchSpec b := ...
\end{lstlisting}

The formulas \code{allFoldSpec}, \code{anyFoldSpec}, and
\code{leastSearchSpec} are built from equations and bounded numerical
relations, not Lean propositions about function values.  Their proofs use
\code{PRA.qfInd}; the logical bounded-quantifier corollaries are derived after
the Boolean equations.

\subsection{Trace-system definitions}

The fixed symbols of a trace are represented directly:

\begin{lstlisting}[style=lean]
structure PRTraceSystem where
  good   : PR 1
  init   : PR 1
  bound  : PR 1
  step   : PR 3      -- (i, previous, u)
  inv    : PR 3      -- (i, state, u)
  accept : PR 2      -- (state, u)

def PRTraceSystem.run (T : PRTraceSystem) : PR 2 :=
  .prec T.init T.step
\end{lstlisting}

Application helpers for arities one through three should be defined using
\code{Fin.cases}.  Define \code{IsOne t := t = succ zero}.  With the de Bruijn
convention documented in the source, define open formulas:

\begin{lstlisting}[style=lean]
def BaseFm  (T : PRTraceSystem) : Fm PTm 1 := ... -- free u
def StepFm  (T : PRTraceSystem) : Fm PTm 3 := ... -- free i,s,u
def FinalFm (T : PRTraceSystem) : Fm PTm 2 := ... -- free s,u
def ResultFm (T : PRTraceSystem) : Fm PTm 0 :=
  Fm.all (... Good(u) -> Accept(run(bound u,u),u) ...)

structure TraceProofs (T : PRTraceSystem) where
  base  : [] |-[PRAAx] BaseFm T
  step  : [] |-[PRAAx] StepFm T
  final : [] |-[PRAAx] FinalFm T
\end{lstlisting}

The intended central theorem is:

\begin{lstlisting}[style=lean]
def traceReflection (T : PRTraceSystem) (h : TraceProofs T) :
    [] |-[PRAAx] ResultFm T := by
  -- instantiate h.base at u
  -- rewrite run(0,u) using PRA.precZero
  -- apply PRA.qfInd to Inv(i, run(i,u), u)
  -- in the step, instantiate h.step and rewrite with PRA.precSucc
  -- specialize the universal invariant at bound(u)
  -- instantiate h.final and cut
  -- discharge Good(u), then universally introduce u
  ...
\end{lstlisting}

This theorem is deliberately a \code{def}, not merely an existence theorem:
it computes the final derivation from the three premise derivations.  The
construction should use only the existing proof constructors, so its axiom
audit is inherited from \code{ConciseIPRA}.

\subsection{Proof-producing finite-program DSL}

After the raw trace theorem works, define a small typed syntax of finite
programs:

\begin{lstlisting}[style=lean]
inductive FProg (input output : Nat) where
  | pure    : PR input -> FProg input output
  | seq     : FProg a b -> FProg b c -> FProg a c
  | loop    : PR input -> PR state -> PR (...) -> FProg input output
  | branch  : PR (...) -> FProg input output ->
              FProg input output -> FProg input output
  | check   : PR (...) -> FProg input input
\end{lstlisting}

The actual indices should track state arities and contracts more carefully than
this sketch.  Compile every program to a \code{PRTraceSystem}.  A
\code{CertifiedProg} pairs syntax with local \code{Prf PRAAx} obligations.
Prove by structural recursion that compilation returns the global soundness
proof.  This automates phase composition and eliminates repeated hand-written
induction scaffolding across theorem packets.

\subsection{Relational finite kernel}

A semantic prototype can use ordinary Lean finite data, but the integrated
module should expose fixed numerical symbols:

\begin{lstlisting}[style=lean]
structure PRLayeredRelation where
  bound  : PR 2       -- (k,p)
  good   : PR 3       -- (a,k,p)
  compat : PR 4       -- (b,a,k,p)

def viable (R : PRLayeredRelation) : PR 4 := ...
def reach (R : PRLayeredRelation) : PR 3 := ...
def fullViable (R : PRLayeredRelation) : PR 4 := ...
\end{lstlisting}

Suggested embedded endpoints are:

\begin{lstlisting}[style=lean]
def viableIffPath (R : PRLayeredRelation) :
  [] |-[PRAAx] viablePathSpec R := ...

def leastViablePathSound (R : PRLayeredRelation) :
  [] |-[PRAAx] leastPathSpec R := ...

def fullViableSplit (R : PRLayeredRelation) :
  [] |-[PRAAx] fullViableSpec R := ...
\end{lstlisting}

Implement viability first as a depth recursion, then add the table algorithm
and prove the two implementations equal.  The depth recursion has the simplest
proof; the table version has the useful execution complexity.

\subsection{Shrinking-viability module}

The module \code{Finitary/Shrinking.lean} should parameterize over coded
rational observations and a coded pseudometric.  All assumptions are
object-language derivations:

\begin{lstlisting}[style=lean]
structure ShrinkingProofs (R : PRLayeredRelation)
    (obs : PR 3) (dist eps delta horizon : ...) where
  nonempty : [] |-[PRAAx] NonemptyCloudFm ...
  monotone : [] |-[PRAAx] HorizonMonotoneFm ...
  diameter : [] |-[PRAAx] CloudDiameterFm ...
  edge     : [] |-[PRAAx] EdgeBoundFm ...
  metric   : MetricProofs dist

def shrinkingViability (h : ShrinkingProofs ...) :
  [] |-[PRAAx] ShrinkingBoundFm ... := ...
\end{lstlisting}

The proof dependencies are:

\begin{enumerate}
\item least bounded search in the full-viable cloud;
\item extraction of a least path certificate from \code{fullViable};
\item truncation of a valid path;
\item finite telescoping of edge distances; and
\item one triangle inequality combining cloud diameter and path motion.
\end{enumerate}

A corollary consumes a tail modulus and constructs a regular-name theorem.
Do not define a quotient real or invoke classical choice.

\subsection{Rational and precision modules}

The modules \code{Rational.lean}, \code{ErrorBudget.lean}, and
\code{RegularName.lean} should prove:

\begin{itemize}
\item cross-multiplication correctness under positive denominators;
\item triangle and product-error inequalities;
\item finite-sum and finite-product accumulators;
\item regularity/extensionality of addition and multiplication;
\item totalized inverse under an apartness guard;
\item overlap branching and exact branching under a gap; and
\item composition of precision ledgers.
\end{itemize}

Most proofs should normalize to polynomial integer inequalities and use small
reusable order lemmas.  A reflection tactic may generate proof terms for
closed numeral calculations, but the universal algebraic lemmas must remain
explicit derivations in the object theory.

\subsection{Application interface}

Each theorem packet should expose a record of the following form:

\begin{lstlisting}[style=lean]
structure PacketTarget where
  params : Nat
  assumptions : List (Fm PTm params)
  conclusion  : Fm PTm params
  derivation  : assumptions |-[PRAAx] conclusion

 def PacketTarget.toNative (T : PacketTarget) :
   NPrf (T.assumptions.map NFm.ofPRA) (NFm.ofPRA T.conclusion) :=
   (pra_iff_native).mp T.derivation
\end{lstlisting}

For an incomplete packet, replace \code{derivation} by an explicitly named
\code{Goal} proposition or a structure of missing producer obligations.  Do
not fill the field with an axiom, \code{opaque} placeholder, or proof of an
ambient semantic surrogate.

\subsection{Dependency-ordered milestones}

A realistic implementation order is:

\begin{enumerate}
\item application helpers and Boolean truth;
\item pairing, bounded lists, signed integers, rationals;
\item folds, bounded minimization, path codes;
\item embedded trace reflection;
\item compiler composition and finite-program DSL;
\item reachability, viability, and least paths;
\item rational metric and shrinking viability;
\item regular names, guards, robust branching, finite nets;
\item exact polynomial/matrix kernel;
\item domain adapters and packet migrations;
\item native export and proof-size measurements.
\end{enumerate}

The first vertical slice should be deliberately small: formalize a finite sum
trace, produce its exact \(\Prf\,\PRAAx\) theorem, and transport it to
\(\NPrf\).  The second slice should be the square-root interval example from
\cref{sec:worked-examples}, because it exercises relational candidates,
viability, shrinking, and rational error bounds.

\subsection{Audit requirements}

For every completed endpoint:

\begin{enumerate}
\item run the repository Lake build;
\item scan source for \code{sorry}, \code{admit}, and new \code{axiom}
      declarations;
\item use \code{\#print axioms} on public theorems;
\item distinguish definitions of goals from proofs of goals;
\item record the pinned ProveIt commit and Lean toolchain;
\item preserve a theorem ledger mapping human claims to Lean names; and
\item retain exact executable checks as regression tests, not as replacements
      for universal proofs.
\end{enumerate}

The accompanying Lean files in this archive are intentionally blueprints: no
Lean executable is available in the build environment, and the generic
\(\Prf\,\PRAAx\) compiler has not been kernel-checked here.  Their purpose is
to make the remaining implementation path explicit and source-accurate.

\section{Worked examples}
\label{sec:worked-examples}

\subsection{A complete finite-sum trace}

Let \(f(p,i)\) be a fixed primitive-recursive integer-valued function and let
\(n(p)\) be a bound.  Define
\[
 S(p,0)=0,
 \qquad
 S(p,i+1)=S(p,i)+f(p,i).
\]
The trace state is the current signed-integer sum, the trace bound is \(n(p)\),
and the invariant at stage \(i\) is the equation
\[
 s=S(p,i).
\]
The base follows from the zero equation for \(S\).  The local step rewrites
both the trace transition and \(S(p,i+1)\) to addition of \(f(p,i)\).  The
final implication identifies the output state with \(S(p,n(p))\).

Although elementary, this is the right first Lean milestone.  It exercises
\code{PR.prec}, \code{PRA.precZero}, \code{PRA.precSucc}, equality transport,
\code{PRA.qfInd}, specialization at a primitive-recursive bound, implication
discharge, and universal introduction.  Once this proof is a reusable
\code{traceReflection} instance, finite products, Horner evaluation, rational
error accumulation, and path generation differ only in their local invariant.

\subsection{Relational square-root approximations}

Consider \(\sqrt2\), but deliberately do not define it as a completed real.
At level \(k\), let the grid denominator be
\(D_k=2^{k+4}\).  A candidate is an integer
\(a\) with \(D_k\le a\le2D_k\), interpreted as
\(q_{k,a}=a/D_k\).  Declare it good when
\begin{equation}
 |q_{k,a}^2-2|\le 2^{-k}.
 \label{eq:sqrt-good}
\end{equation}
Declare adjacent candidates compatible when
\begin{equation}
 |q_{k,a}-q_{k+1,b}|\le 3\cdot2^{-k}.
 \label{eq:sqrt-compat}
\end{equation}
All tests are exact integer inequalities after clearing powers of two.

There are generally several good grid points.  No exact decision about which
one is ``the'' approximation is needed.  Nonemptiness follows from a bounded
integer-square-root search: choose the least \(a\) with \(a^2\ge2D_k^2\), and
compare \(a/D_k\) with the preceding grid point.  The mesh is fine enough to
make one of these points satisfy \eqref{eq:sqrt-good}.

If \(q\) and \(r\) are two good candidates at level \(k\), then
\(q,r\in[1,2]\), and
\[
 |q-r|=\frac{|q^2-r^2|}{q+r}
 \le \frac{|q^2-2|+|r^2-2|}{2}
 \le 2^{-k}.
\]
Thus the viable cloud has diameter at most \(2^{-k}\), independently of which
finite path produced its candidates.  Compatibility gives an edge budget
\(3\cdot2^{-k}\).  A slightly larger precision shift makes the geometric tail
and the horizontal diameter smaller than the regular-name target.  The
shrinking-viability theorem therefore turns the least full-viable grid point
into a primitive-recursive Cauchy name.

The executable companion enumerates these finite grids, computes full-viable
clouds by dynamic programming, and checks the resulting bounds for sample
horizons.  This is regression evidence for the definitions, not a replacement
for the universal iPRA argument above.

\subsection{Roots at a collision}

Consider
\[
 P_t(Y)=Y^2-t^2.
\]
For positive \(t\), the two roots can be labeled by sign.  At \(t=0\), they
collide.  A uniform ordered pair of roots whose first component is always the
negative root and second is the positive root requires an exact sign decision
on represented \(t\) if the presentation permits both signs and zero.

The relational presentation returns two root boxes, counted with
multiplicity, or one cluster of multiplicity two when the separation is below
the requested scale.  Polynomial residuals and Vieta coefficient checks are
finite.  If a certificate supplies \(|t|\ge g>0\), then approximating \(t\)
within \(g/3\) yields its sign, separates the roots by \(2g\), and recovers the
ordered labels.  Thus the same library supports:

\begin{itemize}
\item unconditional multiset/cluster output;
\item an individual approximate root relation;
\item exact ordered roots under a positive separation guard.
\end{itemize}

This pattern extends to eigenvalues, eigenvectors, algebraic branches, and
projective directions.

\subsection{Overlapping coordinate charts}

Suppose a represented point on a curve can be computed in chart \(U_0\) when
\(x<2/3\) and in chart \(U_1\) when \(x>1/3\).  The overlap
\((1/3,2/3)\) is intentional.  Approximate \(x\) within \(1/12\) and compare
the approximation with \(1/2\).  If it is below \(1/2\), use \(U_0\);
otherwise use \(U_1\).  A wrong comparison can occur only when the true point
lies in the overlap, where both formulas are certified.

The branch bit is not an invariant mathematical object, but the output point
is extensional.  This small example captures the preferred treatment of
pivots, local logarithms, rational parametrizations, and analytic charts.

\subsection{Bounded rational points}

Let \(F_p(X,Y)\in\Z[X,Y]\) be a fixed primitive-recursive coefficient family.
A normalized rational pair of height at most \(H\) is coded by integers
\((a,b,c)\) with \(c>0\), bounded absolute values, and the selected gcd/sign
normalization.  The equation
\(F_p(a/c,b/c)=0\) becomes an exact integer equation after multiplying by a
fixed power of \(c\).

The finite enumerator loops over the bounded box, retains normalized solutions,
and removes duplicates by code equality.  Its invariant says that the current
list contains exactly the solutions among the first \(j\) candidate codes.
The trace compiler proves soundness and completeness below \(H\).

A separate height certificate
\[
 F_p(x,y)=0\to h(x,y)\le H
\]
turns this bounded list into the complete list.  Without the certificate, the
same executable search is only a list below a cutoff.  This is the precise
status distinction needed in the Faltings and Siegel settings.

\subsection{A finite-net minimum}

Let \(f:[0,1]^d\to\R\) have a supplied rational modulus.  At requested
precision \(k\), choose a rational grid fine enough that variation inside a
cell is at most \(2^{-k-2}\).  Evaluate \(f\) at the grid points to error
\(2^{-k-2}\), and return the least point with minimum approximate value.
Finite search proves that its displayed value is within \(2^{-k}\) of the
infimum.

If there may be several minimizers, the point itself remains a relation of
near-minimizers.  If a strong-convexity certificate gives
\[
 f(x)\ge f(x_*)+\lambda\|x-x_*\|^2,
\]
then every near-minimizer lies in a rational ball of radius obtained from its
value error.  These balls shrink, so the least full-viable grid point is a
coherent name for \(x_*\).  The value theorem and point theorem are thereby
cleanly separated.

\subsection{End-to-end proof flow}

A mature packet follows the flow
\[
\boxed{\text{codes and guards}}
\longrightarrow
\boxed{\text{finite observations}}
\longrightarrow
\boxed{\text{trace/certificate}}
\longrightarrow
\boxed{\Prf\,\PRAAx}
\longrightarrow
\boxed{\NPrf}.
\]
If outputs are non-unique, insert
\[
\boxed{\text{finite relation}}
\longrightarrow
\boxed{\text{viability}}
\longrightarrow
\boxed{\text{shrinking, if justified}}
\]
between finite observations and the final represented object.  If a deep
certificate producer is unknown, the first four boxes can still be completed
as a conditional theorem whose missing input is explicit.

\section{Retained barriers and impossibility results}
\label{sec:barriers}

A general framework is credible only if it states what encoding cannot repair.
The following barriers explain why the positive interfaces in earlier sections
are deliberately guarded, relational, or certificate-relative.

\subsection{No total universal primitive-recursive evaluator}

\begin{theorem}[Diagonal barrier]
\label{thm:no-universal-pr}
There is no binary primitive-recursive function \(U(e,x)\) such that every
unary primitive-recursive function \(f\) has an index \(e\) with
\(f(x)=U(e,x)\) for all \(x\).
\end{theorem}

\begin{proof}
If such \(U\) existed, then
\[
 g(x)=U(x,x)+1
\]
would be primitive recursive.  By universality, \(g(x)=U(e,x)\) for some
\(e\).  Substitution of \(x=e\) gives
\[
 U(e,e)+1=g(e)=U(e,e),
\]
a contradiction.
\end{proof}

This does not prevent a syntax of primitive-recursive programs with a
partial, clocked, or complexity-bounded interpreter.  It prevents one fixed
total primitive-recursive evaluator from uniformly evaluating every such
program.  Therefore the framework uses fixed evaluator schemas or finite
traces whose individual steps are checked.

\subsection{A cardinality bound is not a search cutoff}

Let \(T(e,t)\) be the primitive-recursive predicate saying that machine \(e\)
halts for the first time at exactly stage \(t\).  For each \(e\), at most one
\(t\) satisfies \(T(e,t)\).  Thus the solution set has the uniform cardinality
bound one.

\begin{theorem}[No uniform cutoff from cardinality]
\label{thm:cardinality-no-cutoff}
There is no total computable function \(H(e)\) such that
\[
 T(e,t)\to t\le H(e)
\]
for every \(e,t\).
\end{theorem}

\begin{proof}
Given such \(H\), simulate machine \(e\) for \(H(e)\) stages.  If it halts,
answer yes.  If it has not halted, it can never halt later, because its first
halting stage would exceed the promised bound.  This decides the halting
problem.
\end{proof}

Consequently, a theorem proving that a curve has at most \(N\) rational points
does not give a height bound, an enumeration stopping rule, or the points
themselves.  Additional arithmetic structure may produce a height bound, but
cardinality alone cannot.

\subsection{Exact sign and LLPO}

Suppose a total operation on arbitrary fast Cauchy names returned one of the
two valid conclusions
\[
 x\le0\qquad\text{or}\qquad x\ge0.
\]
Given a binary sequence \(\alpha\) with at most one nonzero term, form the fast
Cauchy real
\[
 x_\alpha=\sum_{n\ge0}(-1)^n\alpha_n2^{-n-1}.
\]
If the operation returns \(x_\alpha\ge0\), no odd-indexed term can be one; if
it returns \(x_\alpha\le0\), no even-indexed term can be one.  At the zero
sequence both conclusions are valid, so either answer is acceptable.  This is
exactly the lesser limited principle of omniscience (LLPO).

Thus an exact side, sign, orientation, or branch selector at a zero boundary
is not obtained merely by taking more approximations.  A positive gap makes
the finite comparison valid; an overlap makes either branch sound; otherwise
the output should remain relational.

The rotation-axis obstruction is an instance.  Near the identity, the
oriented vector part of a quaternion can approach zero from incompatible
directions.  Exact unit-axis selection would decide a discontinuous side
problem.  Quaternion data with a possibly zero vector part, or approximate
axes, remain stable.

\subsection{Arbitrary Cauchy coefficients do not yield exact Puiseux data}

Consider, for a nonnegative represented real \(a\),
\[
 P_a(X,Y)=Y^2-X^2-a.
\]
Suppose a uniform algorithm returned a ramification exponent \(e>0\), a
finite Laurent lower bound, and exact computable names of coefficients of a
formal root
\[
 y(t)=\sum_j y_jt^j,
 \qquad y(t)^2=t^{2e}+a.
\]
Finite coefficient comparison gives the following dichotomy.

If \(a=0\), all coefficients below degree \(e\) vanish and
\(y_e^2=1\), so \(|y_e|=1\).  If \(a>0\), then \(y_0^2=a\) is apart from zero;
induction through degrees below \(2e\) forces \(y_e=0\).  An approximation to
the algorithm's own coefficient \(y_e\) therefore decides whether \(a=0\).

Now let \(a_m\) have the fixed primitive-recursive fast Cauchy approximation
\[
 A(m,s)=\sum_{k<s}2^{-k-1}H(m,k),
\]
where \(H(m,k)\) says that machine \(m\) has halted by stage \(k\).  If the
machine never halts, \(a_m=0\); if its first halting stage is \(h\), then
\(a_m=2^{-h}>0\).  The hypothetical Puiseux producer would decide halting.

The obstruction concerns exact individually accessible branch coefficients
from arbitrary Cauchy input.  It does not rule out approximate root clusters,
conditional simple branches, or exact Puiseux computation from finite
algebraic presentations supplying equality and splitting data.

\subsection{Metastability does not supply a Cauchy modulus}

There are primitive-recursive increasing bounded rational sequences with no
computable Cauchy modulus.  One standard construction uses stage-bounded
halting information:
\[
 s_n=\sum_{e<n\,:\,e\text{ halts within }n\text{ stages}}2^{-e-1}.
\]
The map \(n\mapsto s_n\) is primitive recursive, rational, increasing, and
bounded by one.  If a computable modulus \(M(r)\) made every later variation
smaller than \(2^{-r}\), then to decide whether machine \(e\) ever halts one
would inspect stage \(M(e+2)\).  A later first contribution of weight
\(2^{-e-1}\) would exceed the allowed future variation; hence absence by that
stage would imply nonhalting.  This is impossible.

Finite stable-window or metastability theorems can still be primitive
recursive and useful.  What fails is the uniform conversion of such windows
into one global convergence rate.  The shrinking-viability theorem therefore
requires a tail modulus as an explicit input.

\subsection{Qualitative uniqueness is not a quantitative selector}

A statement that exactly one mathematical object satisfies a predicate does
not automatically give a primitive-recursive rate at which finite candidate
clouds shrink.  The unique object may be encoded by a computation whose
convergence has no computable modulus.  To construct a name, the framework
requires an isolation radius, strong-convexity constant, residual-to-distance
bound, contraction factor, or equivalent quantitative determinacy data.

This is why the report distinguishes:

\begin{enumerate}
\item uniqueness as a logical proposition;
\item finite near-uniqueness at a requested scale; and
\item a primitive-recursive modulus witnessing how the scale improves.
\end{enumerate}

Only the third supports a uniform coherent-name constructor.

\subsection{A semantic Lean proof does not automatically reflect into iPRA}

Lean can prove an ambient theorem about a recursive function using induction,
quotients, classical choice, algebraic hierarchies, or library facts.  There is
no general reflection principle turning an arbitrary such theorem into a term
of \(\Prf\,\PRAAx\).  A soundness theorem goes in the other direction: an
object-language proof has a semantic interpretation.

The positive solution is deliberately smaller.  For the restricted DSL of
Boolean folds, bounded traces, and certified arithmetic transducers, define a
syntax-directed compiler that constructs the object-language proof from local
object-language obligations.  This is feasible because every proof step and
induction invariant is exposed.  It is not a universal reflection theorem for
Lean.

\subsection{Boundary summary}

The framework does not claim:

\begin{itemize}
\item a universal total primitive-recursive evaluator;
\item exact discontinuous selectors from arbitrary Cauchy names;
\item height cutoffs from cardinality bounds;
\item a Cauchy rate from metastability alone;
\item quantitative isolation from qualitative uniqueness alone; or
\item automatic reflection of arbitrary semantic Lean proofs.
\end{itemize}

Each refused claim has a positive replacement: fixed evaluators or trace
certificates, relational output or positive gaps, explicit height producers,
finite-window theorems or tail moduli, shrinking certificates, and a restricted
proof-producing compiler.

\section{Further research questions and implementation program}
\label{sec:research}

The framework turns a diffuse collection of difficulties into a relatively
small implementation and mathematics agenda.

\subsection{Kernel and compiler questions}

\paragraph{A minimal proof-producing DSL.}
What is the smallest finite-program language whose certified compiler covers
most theorem packets?  The language should support primitive recursion,
bounded folds, adaptive queries, finite search, guarded branching, and phase
composition, while producing compact \(\Prf\,\PRAAx\) terms.  It would be
valuable to compare direct proofs in ordinary iPRA with proofs compiled first
to the native scalar or clause calculi.

\paragraph{Proof-size control.}
Na\"ively expanding Boolean folds and rational arithmetic may create enormous
proof terms.  Which derived rules, sharing mechanisms, or certificate
languages minimize proof size without enlarging the trusted kernel?  Can a
small native certificate be checked by a proved object-language verifier and
then reflected to the ordinary presentation?

\paragraph{Verified code generation.}
Can the Python-style exact checkers already present in many packets be replaced
by a restricted source language with a verified compiler to \(\PR\) syntax,
trace invariants, and proof obligations?  The trusted output should remain a
Lean derivation, not the external compiler's assertion of correctness.

\subsection{Relational mathematics questions}

\paragraph{General finite diagrams.}
The report treats layered paths.  Many applications naturally form finite
DAGs, trees, or hypergraphs of compatibility constraints.  Develop a generic
bounded CSP library and identify the weakest local consistency conditions that
permit efficient proof-producing amalgamation.

\paragraph{Quantitative determinacy extractors.}
Build reusable theorems that derive viable-cloud diameter bounds from:
strong convexity, inverse-function estimates, algebraic separation, Rouch\'e
bounds, contraction mappings, monotone brackets, and projective spectral gaps.
This would make shrinking viability a practical bridge rather than a manual
argument in each packet.

\paragraph{Multiset and cluster representations.}
Design one canonical finite representation for approximate multisets with
multiplicity, Hausdorff matching, and cluster refinement.  Prove that
polynomial root sets, eigenvalue sets, and finite zero sets are extensional in
this representation, including collision points.

\subsection{Analysis library questions}

\paragraph{A common error-budget algebra.}
Formalize rational error expressions, monotone precision transformers, and
automatic composition.  Can a small verified normalizer discharge the routine
linear and polynomial inequalities generated by finite sums, products,
quadrature, and series tails?

\paragraph{Contour and zero-count certificates.}
Develop a finite polygonal-contour library with boundary apartness, winding,
argument increments, Cauchy estimates, and Rouch\'e certificates.  This would
connect the zeta, residue, FTA, and analytic-branch packets through one kernel.

\paragraph{Represented measures by finite tests.}
Extend FRPs so that challenges are bounded rational test functions and answers
are approximate integrals.  Determine which probability results need only
finite laws plus approximation moduli and which benefit from the more general
interface.

\subsection{Number-theoretic and geometric producer questions}

\paragraph{Height-bound producers.}
For each constructive special case of Siegel, Faltings, or Mordell--Weil in the
repository, isolate the exact numerical theorem that produces a height bound or
covering representatives.  Can power-gap, Runge, descent, and local-height
arguments be factored into a shared producer library?

\paragraph{Effective analytic constants.}
For prime-distribution packets, specify machine-checkable constant tables and
then formalize the strongest currently effective producer results.  Where a
classical theorem is ineffective, record the verifier theorem without claiming
a producer.

\paragraph{Projective and quaternionic geometry.}
Develop projective points, lines, eigenvectors, axes, and quaternions as
guarded relational presentations.  Quantify precisely which classical exact
geometry statements require LLPO-like selectors and which admit stable
projective or multivalued replacements.

\subsection{Repository migration questions}

\paragraph{Packet triage.}
For each of the 58 delivered packets, fill a uniform schema containing its
codes, finite observations, guards, precision ledger, trace system,
object-language target, and missing producers.  This would replace prose-only
status with a machine-readable dependency graph.

\paragraph{Reference vertical slices.}
Three migrations would exercise most of the architecture:

\begin{enumerate}
\item an exact finite identity package, such as a divisor or trigonometric
      polynomial identity;
\item an analytic approximation package, such as Stirling or a finite
      Gamma-product theorem; and
\item a relational package, such as polynomial root clusters or rotation
      quaternions/axes.
\end{enumerate}

Each should end in a kernel-checked \(\Prf\,\PRAAx\) theorem and native export.

\paragraph{Status automation.}
Can continuous integration verify that every public claim has a corresponding
Lean declaration, that all intended declarations build, that no placeholder
axioms are introduced, that PDFs rebuild, and that executable examples agree
with the formal definitions?  A theorem ledger should fail closed: an absent
proof remains visibly absent.

\subsection{Conceptual question}

The broadest question is whether a substantial portion of effective
mathematics admits a small common language of finite observations, relational
witnesses, quantitative determinacy, and proof-producing traces.  The present
report gives evidence for an affirmative answer, while retaining precise
boundaries where deep producer theorems or constructive obstructions remain.
Completing the compiler and two or three representative migrations would test
that claim much more decisively than adding further prose-only theorem packets.

\appendix
\section{Exact theorem schemas and packet checklist}
\label{app:schemas}

This appendix collects the core theorem shapes in a form suitable for direct
translation into Lean declarations.

\subsection{Boolean bounded-formula compiler}

For every bounded formula \(\theta(\bar x)\) formed from equations between
fixed primitive-recursive terms, propositional connectives, and bounded
quantifiers, construct a primitive-recursive Boolean
\(b_\theta(\bar x)\) and derivations
\[
 []\vdash[\PRAAx]
 \bigl(b_\theta(\bar x)=1\to\theta(\bar x)\bigr),
 \qquad
 []\vdash[\PRAAx]
 \bigl(\theta(\bar x)\to b_\theta(\bar x)=1\bigr).
\]
The construction is structural recursion on \(\theta\).  The bounded
quantifier cases use the fold-correctness derivations.

In the compiled Lean implementation, the source syntax should be a restricted
inductive type carrying its free-variable arity.  Compilation returns the
Boolean \(\PR\) symbol, the object-language formula, and both derivations.

\subsection{Proof-producing trace reflection}

Given fixed symbols \(T=(g,i,b,s,v,a)\), where \(g\) is the good-input bit,
\(i\) the initial state, \(b\) the bound, \(s\) the step, \(v\) the invariant,
and \(a\) the acceptance bit, let \(r=\code{PR.prec }i\ s\).  The input record
contains derivations of:
\begin{align*}
 g(u)=1&\to v(0,i(u),u)=1,\\
 g(u)=1\wedge v(j,z,u)=1&\to v(j+1,s(j,z,u),u)=1,\\
 g(u)=1\wedge v(b(u),z,u)=1&\to a(z,u)=1.
\end{align*}
The output is a derivation of
\[
 \forall u\;\bigl(g(u)=1\to a(r(b(u),u),u)=1\bigr).
\]
No semantic premise about the standard natural numbers appears in the theorem
signature; every premise is already an object-language derivation.

\subsection{Finite path compiler}

For fixed \(B,\Good,\Compat\), define the path checker
\(\mathsf{PathOK}(p,k,d,a,c)\).  Its fields assert:

\begin{enumerate}
\item the code \(c\) has length \(d+1\);
\item its first entry is \(a\);
\item every entry is below the candidate bound and good; and
\item every adjacent pair is compatible.
\end{enumerate}

The viability theorem is the pair of derivations
\begin{align*}
 \Viable(p,k,d,a)=1
 &\to
 \exists c<C(p,k,d)\;\mathsf{PathOK}(p,k,d,a,c)=1,\\
 \mathsf{PathOK}(p,k,d,a,c)=1
 &\to \Viable(p,k,d,a)=1,
\end{align*}
where \(C\) is a displayed primitive-recursive code bound.  The existential
witness in the first formula can be the least path code or the recursively
constructed least viable path.

\subsection{Shrinking viability}

The exact numerical conclusion is
\[
 \forall k\le\ell\quad
 d(\Obs(p,k,s(p,k)),\Obs(p,\ell,s(p,\ell)))
 \le\varepsilon(p,k)+
       \sum_{j=k}^{\ell-1}\delta(p,j).
\]
The assumptions are:

\begin{enumerate}
\item \(s(p,k)\) is the least member of the nonempty cloud
      \(\FullV(p,H(p,k),k,-)\);
\item \(H\) is nondecreasing and dominates its level;
\item the cloud diameter is at most \(\varepsilon(p,k)\);
\item compatible edges move observations by at most \(\delta(p,k)\); and
\item the coded distance satisfies symmetry and triangle inequality on the
      guarded domain used in the theorem.
\end{enumerate}

The regular-name corollary takes a separate tail modulus.  The theorem does not
infer summability of \(\delta\) or convergence of \(\varepsilon\).

\subsection{Finite-net transfer}

A net certificate has a finite point list \(N\), a coverage checker, and a
radius \(\rho\).  A modulus certificate has a bound \(\eta\) at radius
\(\rho\).  The embedded target is:
\[
 \mathsf{NetOK}(p,N,\rho)=1\wedge
 \mathsf{ModulusOK}(p,\rho,\eta)=1\wedge
 \All_{y\in N}[f(y)\ge L]=1
 \to
 \mathsf{GlobalLower}(p,f,L-\eta)=1.
\]
The meaning of \(\mathsf{GlobalLower}\) depends on the represented domain; its
proof must reduce an arbitrary challenge to one finite coverage witness.

\subsection{Height bound and list}

For a normalized point-code predicate \(P(p,x)\) and height \(h(x)\), define:
\begin{align*}
 \mathsf{BoundOK}(p,H)&:=\forall x\,[P(p,x)\to h(x)\le H],\\
 \mathsf{ListOK}(p,L)&:=
   \bigl(\forall i<|L|\;P(p,L_i)\bigr)\wedge
   \bigl(\forall x\;P(p,x)\to\exists i<|L|\;x=L_i\bigr).
\end{align*}
The two constructors are:
\begin{align*}
 \mathsf{enumerate}(p,H)&:\mathsf{BoundOK}(p,H)	o
     \mathsf{ListOK}(p,\mathsf{enumerate}(p,H)),\\
 \mathsf{maxHeight}(L)&:\mathsf{ListOK}(p,L)	o
     \mathsf{BoundOK}(p,\mathsf{maxHeight}(L)).
\end{align*}
The first depends on a primitive-recursive finite bound on normalized codes of
height at most \(H\), which must be stated for the chosen coding.

\subsection{Uniform packet checklist}

Before a theorem packet is called embedded, record all fields below.

\begin{longtable}{@{}p{0.28\textwidth}p{0.64\textwidth}@{}}
\toprule
Field & Required evidence\\
\midrule
\endhead
Pinned source & ProveIt commit and exact imported module paths.\\
Statement class & Exact, approximate, certificate-relative, relational, or
obstruction theorem.\\
Parameter codes & Natural-number coding and domain guard.\\
Fixed evaluators & Exact \(\PR\) symbols; no unmentioned program input.\\
Challenges/answers & Bounds, verifier, and observational meaning.\\
Compatibility & Adjacent or finite-diagram relation, if outputs are relational.\\
Precision ledger & All input queries, shifts, truncations, margins, and errors.\\
Trace & Initial state, step, bound, invariant, and acceptance predicate.\\
Local proofs & Actual \(\Prf\,\PRAAx\) derivations of base, step, and final
contracts.\\
Search data & Net, height bound, isolating region, cover, or finite candidate
bound where required.\\
Producer status & Proved, conditional, externally certified, or open.\\
Object endpoint & Exact \code{Fm PTm} conclusion and context.\\
Native export & Lean theorem applying \code{pra\_iff\_native} to the completed
endpoint.\\
Audit & Build log, axiom report, placeholder scan, and theorem ledger.\\
Executable checks & Exact examples and negative tests, explicitly labeled as
regression evidence.\\
\bottomrule
\end{longtable}

A packet failing one field may still contain valuable mathematics.  The
checklist identifies the missing layer without conflating it with the layers
already complete.

\section{Repository source audit and inheritance}
\label{app:source-audit}

\subsection{Snapshot and scope}

The repository was inspected read-only at main commit
\[
 \code{b88e843c31b50dfe461bfbd0380166ca422cd3ab}.
\]
The audit focused on the checked iPRA backend, the proof-packet index, all
foundation-packet summaries, and representative obstruction reports.  This is
not a verification of all theorem packets.  Claims inherited from earlier
packets are credited below; the present report's additions are the unified
proof-producing compiler interface and the shrinking-viability synthesis.

\subsection{Primary checked backend}

\begin{longtable}{@{}p{0.38\textwidth}p{0.54\textwidth}@{}}
\toprule
Repository path & Use in this report\\
\midrule
\endhead
\path{Logic/PrimitiveRecursiveArithmetic/README.md} & Status boundary between
checked \code{ConciseIPRA} and delivered research packets; packet count and
absence of completed packet derivations.\\
\addlinespace
\path{Logic/PrimitiveRecursiveArithmetic/ConciseIPRA/README.md} & Native scalar
calculus, ordinary iPRA presentation, and formalization overview.\\
\addlinespace
\path{ConciseIPRA/Lean/ConciseIPRA/PR.lean} & Exact definitions of
\code{PR}, \code{PTm}, \code{PRAAx}, the defining-equation lemmas, and
\code{PRA.qfInd}.\\
\addlinespace
\path{ConciseIPRA/Lean/ConciseIPRA/Deduction.lean} & Exact derivation type
\code{Prf}, substitution, weakening, cut, quantifier rules, and equality
transport.\\
\addlinespace
\path{ConciseIPRA/Lean/ConciseIPRA/NativeBiInterpretation.lean} & Exact
statement of \code{pra\_iff\_native} and the native/ordinary translations.\\
\bottomrule
\end{longtable}

\subsection{Packet catalog and foundation reports}

\begin{longtable}{@{}p{0.42\textwidth}p{0.50\textwidth}@{}}
\toprule
Repository path & Inherited contribution or status\\
\midrule
\endhead
\path{docs/proofs/README.md} & Catalog of 58 packets, domain grouping, and
packet-specific status summaries.\\
\addlinespace
\path{foundations/ipra_embedding_framework/} & Finite observations,
bounded Boolean predicates, local analytic branches, stability windows, and
finite algebra certificates.\\
\addlinespace
\path{foundations/ipra_embedding_framework_2026-11-15/} & Certified
Finitary Presentations, finite-trace reflection at the human-proof level,
precision transformers, guards, finite nets, and an explicit ledger saying the
generic embedded \code{Prf PRAAx} proof is not yet implemented.\\
\addlinespace
\path{foundations/ipra_evidence_framework/} & Relational finite outputs,
compatibility, and backward finite dynamic programming for witness
amalgamation.\\
\addlinespace
\path{foundations/ipra_finite_observation/} & Uniform finite-table compiler,
degree-uniform finite constructions, quantitative minimizers, and checked
ambient finite-table prototypes.\\
\addlinespace
\path{foundations/ipra_framework/} & Finite contracts, guarded induction,
regular-name arithmetic, stable windows, local branches, and compiled embedded
examples.\\
\addlinespace
\path{foundations/ipra_framework_2026-11-15/} & Bounded search/descent,
certificate composition, quantitative names, and guarded inverse.\\
\addlinespace
\path{foundations/ipra_uniform_certificates/} & Finite-net lifting, root
names under sign/separation certificates, and precision-contract composition.\\
\bottomrule
\end{longtable}

\subsection{Representative difficulty reports}

\begin{longtable}{@{}p{0.42\textwidth}p{0.50\textwidth}@{}}
\toprule
Repository path & Boundary used here\\
\midrule
\endhead
\path{geometry/euler_rodrigues_ipra/README.md} & LLPO obstruction to a
uniform exact unit axis and positive replacements using quaternions,
approximation, or a separation gap.\\
\addlinespace
\path{algebra/puiseux_ipra/evidence/scope_obstruction.md} & Halting reduction
against uniform exact Puiseux coefficient production from arbitrary Cauchy
coefficient names.\\
\addlinespace
\path{diophantine-number-theory/faltings_ipra/README.md} & Distinction between
finiteness/cardinality and a height bound; complete-list--bound equivalence;
status of general Faltings.\\
\addlinespace
\path{foundations/ipra_embedding_framework_2026-11-15/metadata/difficulty_matrix.csv} & Cross-domain list of reported obstacles and residual
status.\\
\bottomrule
\end{longtable}

\subsection{What is new in this report}

The following synthesis is not copied as a theorem from one prior packet:

\begin{enumerate}
\item a precise proof-producing interface in which the base, step, and final
      trace obligations are themselves \(\Prf\,\PRAAx\) terms and the generic
      compiler is a Lean function returning the completed derivation;
\item the integration of relational finite outputs with two-sided
      full-horizon viability;
\item the shrinking-viability theorem, which chooses the least full-viable
      candidate at changing finite horizons and proves coherence from cloud
      diameter and edge-tail bounds, without selecting one infinite branch;
\item a unified migration boundary separating finite kernel, certificate
      producer, object-language endpoint, and native export; and
\item a dependency-ordered implementation program against the exact current
      \code{ConciseIPRA} API.
\end{enumerate}

The report makes no claim that these additions have already been kernel-checked
in Lean.  The human proofs are intended to guide that implementation, and the
accompanying exact Python examples validate finite definitions only.


\end{document}
